\documentclass[11pt,a4paper]{article}

\usepackage[utf8]{inputenc}
\usepackage[T1]{fontenc}
\usepackage{lmodern}
\usepackage{amsmath, amssymb, amsthm}
\usepackage{mathtools}
\usepackage{geometry}
\usepackage{hyperref}
\usepackage{xcolor}
\usepackage{booktabs}
\usepackage{enumitem}
\usepackage{microtype}
\usepackage{listings}
\usepackage{verbatim}
\usepackage{seqsplit}

\geometry{margin=2.8cm}

% Help LaTeX break long monospace identifiers (camelCase + underscores) to
% avoid overfull hboxes in v3 sections that reference Lean declaration names
% and Phase-7 CSV file names.
\setlength{\emergencystretch}{3em}
\sloppy
% Breakable typewriter macro: like \texttt but allows line breaks at every
% character.  Use for identifiers that overflow the line, not as a default.
\newcommand{\ttbr}[1]{\texttt{\seqsplit{#1}}}
\hypersetup{
  colorlinks=true,
  linkcolor=blue!50!black,
  citecolor=blue!50!black,
  urlcolor=blue!50!black,
  pdftitle={A Spectral Reduction of the Collatz Conjecture via Phantom Orbit Shadowing (v3)},
  pdfauthor={Piero Borgatta}
}

\newtheorem{theorem}{Theorem}
\newtheorem{lemma}[theorem]{Lemma}
\newtheorem{proposition}[theorem]{Proposition}
\newtheorem{corollary}[theorem]{Corollary}
\theoremstyle{definition}
\newtheorem{definition}[theorem]{Definition}
\newtheorem{conjecture}[theorem]{Conjecture}
\newtheorem{remark}[theorem]{Remark}
\newtheorem{example}[theorem]{Example}

\DeclareMathOperator{\nutwo}{\nu_2}
\newcommand{\Z}{\mathbb{Z}}
\newcommand{\Zp}{\mathbb{Z}_2}
\newcommand{\Q}{\mathbb{Q}}
\newcommand{\N}{\mathbb{N}}
\newcommand{\R}{\mathbb{R}}
\newcommand{\full}{\mathrm{FULL}}
\newcommand{\core}{\mathrm{CORE}}
\newcommand{\tail}{\mathrm{TAIL}}
\newcommand{\spec}{\rho}

\title{
  A Spectral Reduction of the Collatz Conjecture \\
  via Phantom Orbit Shadowing
  \\[4pt]\large\normalfont (Version 3 -- phantom-set taxonomy and Lean-checked finite operator certificates)
}
\author{
  Piero Borgatta\thanks{Email: \texttt{info@pieroborgatta.com}.\
    ORCID: \href{https://orcid.org/0009-0001-8025-2405}{0009-0001-8025-2405}.}
}
\date{May 2026 (v3) \\[2pt]
  \small\normalfont
  Source: \href{https://github.com/PieroBorgatta/Collatz}{\texttt{github.com/PieroBorgatta/Collatz}}
  \\[2pt]
  \small\normalfont
  Version 1: \href{https://doi.org/10.5281/zenodo.20021538}{\texttt{10.5281/zenodo.20021538}} \quad
  Version 2: \href{https://doi.org/10.5281/zenodo.20098868}{\texttt{10.5281/zenodo.20098868}} \quad
  Version 3: \href{https://doi.org/10.5281/zenodo.20160154}{\texttt{10.5281/zenodo.20160154}}
}

\begin{document}
\maketitle

\begin{abstract}
We propose a structural reformulation of the Collatz conjecture in
terms of a finite-dimensional spectral problem. Starting from the
empirical observation that ``rebel'' orbits of the Syracuse map
shadow, for finite stretches, periodic rational points of the
$2$-adic dynamics associated to negative rational fixed points
(\emph{phantom cycles} in the terminology of \cite{chang2026},
introduced concurrently and independently of this work; the
underlying object is classical and goes back to
\cite{lagarias1985}), we construct an episode graph on the space of
phantom orbits. The strongly connected component (SCC) of this graph
that supports the slowest descent is encoded as a finite weighted
transfer operator on a refined phase state
$(\nu_2(t), \mathrm{odd}(t)\bmod 4, h\bmod 4)$ at the critical node
$(k=12,\, c=2,\, b=1)$.

We prove an exact congruential shadowing lemma reducing the question
of arbitrarily long shadowing episodes to a residue condition
$n\equiv q_w\pmod{2^{bA+1}}$, where $q_w$ is the $2$-adic
representative of a phantom word $w$ and $A=|w|_{\nu_2}$ its period
sum. The lemma is the iterated step-by-step form of the per-block
$2$-adic repulsion identity of \cite[Prop.~7.4]{chang2026} and is
equivalent in content to the classical inverse formula for parity
vectors \cite[Lemma~1]{rozier2025}; the iterated form is what feeds
the episode graph construction. \textbf{The lemma and its
no-infinite-shadowing corollary have been mechanically verified in
Lean~4 with Mathlib} (Section~\ref{sec:formalization}); the project
is \texttt{sorry}-free.

We then build, for each truncation depth $T$ and high-bit lift count
$j$, a finite matrix $\full_{T,j}$ encoding first-return behavior at
the critical node, decompose it as
$\full_{T,j} = \core_{T,j} + \tail_{T,j}$, and apply a weighted
Collatz--Wielandt bound \cite{lemmensnussbaum} with respect to the
right Perron eigenvector $v$ of $\full_{T,j}$. Our numerical results
show
\[
\spec(\full_{T,j})
\le
\max_i (\core_{T,j}\, v)_i / v_i + \max_i (\tail_{T,j}\, v)_i / v_i
\le 0.052
\quad
\text{for $T \le 16$ and $j \le 32$,}
\]
with monotone but geometrically decaying increments in both $T$ and
$j$, and we extend the bound to a cross-node operator on the SCC
obtaining $\spec(M_{\mathrm{cross}}) \le 0.0366$. The cross-node
spectrum is \emph{lower} than the single-node bound at the critical
vertex.

\textbf{New in v3.} (i) A finite phantom-set taxonomy at depth
$K_0 = 16$ via M\"obius inversion of necklace counts: all $1247$
primitive expansive compositions with $K \le 16$ are enumerated, their
2-adic representatives computed in exact arithmetic, and the orbit
harness identifies a single nontrivial taxonomy SCC of $1222$ raw
states (Section~\ref{sec:taxonomy}). (ii) A \emph{deterministic} finite
residue-cell Collatz--Wielandt certificate on the compressed $(K,b)$
macro-state space (37 states): exact rational max ratio
$90833233962213/129559208330288 < 3/4$, with sensitivity at
$\mathrm{lift\_bits} \in \{4,5,6\}$ confirming the bound under
finer residue refinement. \textbf{The certificate is imported into
Lean as \ttbr{k16s16KDeterministicGeneratedSpectralRadiusBound},
which proves the bound on Mathlib's real \texttt{spectralRadius}.}
(iii) A Lean-checked finite-operator pipeline for the original critical
node: an exact $T = 10$, $j = 32$ majority-signature core/tail
decomposition is imported as a verified
\texttt{OperatorDecomposition}, the weighted Collatz--Wielandt bound
is proved generically in Lean (\texttt{Bound.lean}), and the
single-node spectral-radius bound
$\spec(\full_{10,32}) \le 97/2000 = 0.0485$ is exposed as
\ttbr{t10j32HighBitTailSpectralRadiusBound\_97\_2000}, with no
\texttt{axiom}, \texttt{sorry}, or \texttt{admit} in the generated
certificate (Section~\ref{sec:formalization}).

We do not claim a proof of the Collatz conjecture. We claim that the
structural problem at the critical SCC is reduced to two explicit a
priori estimates on finite matrices: the asymptotic boundedness of
the weighted bound under refinement in $T$, and the closure of the
empirical signature support under refinement in $j$. The same
distributional-to-pointwise barrier identified in \cite{chang2026}
manifests itself here in finite-matrix form
(Section~\ref{sec:reduction}). We document supporting numerical
evidence over a five-decade range of parameters, expose the open
estimates as conjectures, and provide all computational scripts and
the Lean 4 formalization as supplementary material.
\end{abstract}

\noindent\textit{Research program note.} This paper is the third
versioned artifact of a personal research program by the author, an
independent researcher, exploring whether iterative collaboration with
large language model assistants can support non-standard attempts at
open mathematical problems. \emph{Verifying the practical limits of
such collaboration is an explicit secondary goal of this work, on
equal footing with making progress on Collatz.} v3 extends v2 along
two axes: the phantom-set taxonomy at fixed depth $K_0$ closes a gap
in the empirical SCC framing of v2 by enumerating the full primitive
universe and certifying a deterministic finite-residue-cell
Collatz--Wielandt bound; and the Lean 4 + Mathlib formalization
is extended beyond the shadowing core to cover the finite episode
graph, the truncated transfer operator with empirical signature
decomposition, the generic weighted Collatz--Wielandt bound, and two
generated finite numerical certificates (the deterministic
phantom-taxonomy certificate at $K_0 = 16$ and the single-node
$T = 10$, $j = 32$ certificate). The methodology, the contribution
split between human and AI, the cross-AI verification strategy, and
the v2$\to$v3 phase log are documented in
Section~\ref{sec:methodology} and in \texttt{METHODOLOGY.md} of the
supplementary archive.
\medskip

\tableofcontents

\section{Introduction}
\label{sec:intro}

The Collatz conjecture asserts that the iteration
\[
C(n) =
\begin{cases}
n/2 & \text{if } n \text{ even,} \\
3n+1 & \text{if } n \text{ odd,}
\end{cases}
\]
applied to any positive integer $n$ eventually reaches the trivial
cycle $1 \to 4 \to 2 \to 1$. Despite eighty years of attention from
combinatorialists, number theorists, and dynamicists, the conjecture
remains open. The strongest unconditional progress is due to Tao
\cite{tao2019}, who established that almost all orbits attain
arbitrarily small values relative to the initial point, in a precise
density sense.

The bottleneck for a deterministic argument is well known: certain
``rebel'' starting values produce orbits that climb to many orders
of magnitude above their initial value before eventually descending.
Empirically these rebels are not isolated curiosities but appear to
form structured families which converge in inverse-tree language to
common dissipative trunks.

This work proposes that the structural mechanism behind such families
is \emph{$2$-adic shadowing}: rebel positive integers follow, for
finite stretches, the trajectory of negative rational fixed points
of the Syracuse map under composition. These ``phantoms'' have no
positive integer realization but locally constrain the $\nu_2$
sequence of nearby positive orbits, producing the observed expansive
corridors. The same picture has been investigated independently and
contemporaneously by Chang \cite{chang2026}, who calls the same
objects \emph{phantom cycles} and develops a complementary
broad-spectrum analysis; we discuss the relationship in detail in
Section~\ref{sec:related} and again in
Section~\ref{sec:comparison}.

We make this framework precise, codify the dynamics as a finite
weighted transfer operator on a refined phase quotient, and prove
numerically a uniform spectral upper bound which reduces the
boundedness of orbits at the critical SCC to an explicit a priori
estimate on a family of finite matrices. The fundamental shadowing
lemma is mechanically verified in Lean~4 with Mathlib
(Section~\ref{sec:formalization}). In v3 we additionally close the
phantom-set scope at depth $K_0 = 16$ via M\"obius enumeration of
primitive cyclic compositions, derive an exact deterministic
finite-residue-cell Collatz--Wielandt certificate on the compressed
$(K, b)$ macro-state space (Section~\ref{sec:taxonomy}), and import
both this certificate and a single-node finite-operator certificate
at $T = 10, j = 32$ into the Lean formalization
(Section~\ref{sec:formalization}).

\subsection{Related work}
\label{sec:related}

The literature on the $3x+1$ problem is vast; we restrict the
discussion to the threads most directly relevant to the framework
developed here. Comprehensive bibliographies of earlier work are
available in \cite{lagarias2010,laga2003bib,laga2009bib}.
The reconnaissance corpus carried over from v2 has a cutoff of early
2026; no new arXiv work in the directly overlapping niche
(\emph{phantom orbit shadowing on the $3x+1$ map}) appeared between
v2 (deposited 2026-05-10) and v3 (this revision, 2026-05-13), so the
related-work tour below is unchanged from v2 in substantive content.
A standing monthly arXiv scan in the relevant subject classes is
described in Section~\ref{sec:methodology-future-search}.

\paragraph{Parity vectors and congruence theory.}
The systematic study of parity vectors of the $3x+1$ map originates
with Everett, Terras, and Lagarias in the late 1970s
\cite{lagarias1985}. The classical inverse formula --- expressing
the residue class of integers realizing a given finite parity
sequence --- appears in this line and is reformulated explicitly
in Rozier \cite[Lemma 1]{rozier2025} as
\[
n \equiv -\sum_{k=0}^{j-1} s_k\, 2^k\, 3^{-\sigma_k(S)}
\pmod{2^j},
\]
for a parity sequence $S = (s_0, \dots, s_{j-1})$ with cumulative
$1$-counts $\sigma_k(S)$. Lemma~\ref{lem:shadowing} below is
mathematically equivalent to a particular packaging of this formula
in terms of the rational fixed point $q_w$ of the period word; we
favor the latter packaging because it is what feeds the episode
graph construction in Section~\ref{sec:episodes} and is the form
naturally formalized in Lean (Section~\ref{sec:formalization}).

\paragraph{Non-Archimedean spectral theory.}
Siegel \cite{siegel2023a,siegel2023b,siegel2026hydra} develops a
non-Archimedean spectral theory of the Collatz map via the
\emph{Numen} function $\chi_q\colon \Zp \to \Z_q$ and a
Correspondence Principle linking $2$-adic values to periodic and
divergent orbits of the accelerated Syracuse map. While the
underlying $2$-adic fixed-point equation
$(2^A - 3^L)\, q_w = C_w$ used here for the phantom representatives
coincides with the equation underlying the periodic-point criterion
of Bohm--Sontacchi exploited in \cite{siegel2023a}, the present
work's focus on finite transfer operators on a refined phase
quotient differs sharply from Siegel's complex-analytic strategy
based on contour integrals and meromorphic continuations of
Dirichlet series.

\paragraph{Phantom cycles and transfer operators (concurrent work).}
The work most closely related to the present one is Chang
\cite{chang2026} (\emph{Exploring Collatz Dynamics with Human-LLM
Collaboration}), an independently developed Human--LLM
collaboration produced through what the author terms a 29-paradigm
program comprising several hundred formal results. Three
overlapping themes are noteworthy:
\begin{enumerate}
  \item Chang \cite[Def.~7.2]{chang2026} introduces \emph{phantom
        cycles} of a signature $\sigma = (k_1, \ldots, k_\ell)$ as
        the unique $2$-adic fixed point of the corresponding block
        map $F_\sigma$, with explicit equation
        $(2^K - 3^\ell)\, \rho = C_\sigma$. This is exactly our
        $q_w$ construction (Section~\ref{sec:phantom-words}) up to
        notation $(\sigma \leftrightarrow w,\ K \leftrightarrow A,
        \ \ell \leftrightarrow L,\ \rho \leftrightarrow q_w)$.
  \item Chang's \emph{$2$-adic repulsion} identity
        \cite[Prop.~7.4]{chang2026},
        $\nu_2(F_\sigma(x) - \rho) = \nu_2(x - \rho) - K$,
        is the per-block form of the iterative
        Lemma~\ref{lem:shadowing} below. Chang's no-prolonged-trapping
        statement \cite[Lemma~9.20]{chang2026} corresponds to
        Corollary~\ref{cor:no-infinite}.
  \item Chang \cite[Theorem 7.15]{chang2026} (\emph{Phantom
        Universality}) establishes that \emph{every} primitive
        cyclic composition of valuations is a phantom family, a
        uniform statement strictly stronger than the empirical
        enumeration we use (Appendix~\ref{app:phantoms}, $k \le 24$).
        The number-theoretic gain rate $R \le 0.0893$ obtained in
        \cite[Theorem~7.19]{chang2026} is conceptually similar to
        --- but quantitatively distinct from, and computed on a
        different object than --- our spectral bounds; see
        Section~\ref{sec:comparison}.
\end{enumerate}
The two works share the core 2-adic phenomenology but differ in
strategy: Chang aggregates over all phantom families of growing
depth, while the present work concentrates on the single critical
SCC with finer phase resolution. Chang's transfer operator $T_K$
on $(\Z/2^K\Z)_{\mathrm{odd}}$ is nilpotent
\cite[Remark~9.17]{chang2026}, whereas our $\full_{T,j}$ on the
refined phase state $(\nu_2(t), \mathrm{odd}(t) \bmod 4,
h \bmod 4)$ has $\spec \approx 0.03$. The fact that the same
underlying barrier --- a distributional-to-pointwise gap --- arises
in both formulations independently is, in our view, evidence that
this barrier is genuine and not an artifact of either framing.

\paragraph{Operator-theoretic approaches.}
Berg and Meinardus \cite{berg1996} introduced a linear operator
$F$ on the Hardy space $H^2(D)$ and established functional
equations for the Collatz iteration. Neklyudov
\cite{neklyudov2022} studies the adjoint operator $T$ on the
Bergman space, proving hypercyclicity and connecting the index of
$\mathrm{Id}-T$ to the number of cycles. Leventides--Poulios
\cite{leventides2020} employ Koopman operators on $\ell^p$ spaces
and $C^*$-algebras. Mori \cite{mori2025} formulates the Collatz
conjecture in terms of irreducibility of associated $C^*$-algebras.
Each of these works lifts the discrete Collatz dynamics to an
infinite-dimensional linear setting; the present work instead
constructs finite-dimensional non-negative matrices from empirical
transition data at the critical phantom node.

\paragraph{Spectral methods for transfer operators.}
The Collatz--Wielandt characterization of the spectral radius of
order-preserving homogeneous maps on cones, due to Lemmens,
Nussbaum, and collaborators \cite{lemmensnussbaum,akiangaubert},
underlies the weighted bound used in
Section~\ref{sec:bound}. Keller--Liverani perturbation theory
\cite{kellerliverani} and Baladi's transfer-operator framework
\cite{baladi} are the natural analytic targets for a possible
proof of Conjecture~\ref{conj:T} (see
Section~\ref{sec:limits}). Recent work on certified spectral
approximation of transfer operators \cite{certifiedspectral} and on
quasi-compactness of Frobenius--Perron operators
\cite{quasicompact} is also relevant infrastructure for that
direction.

\paragraph{$p$-adic dynamics and shadowing.}
Anashin \cite{anashin} surveys the non-Archimedean theory of
discrete systems including Collatz-like maps. Bastos, Caprio, and
Messaoudi \cite{bcm2020} develop a general theory of shadowing
and structural stability for $p$-adic dynamics, not specialized to
Collatz. Laarhoven and de Weger \cite{laarhoven2012} construct
modular Collatz graphs on residue classes $\Z / m \Z$ for
$m = 2^K$, identifying an isomorphism with binary De~Bruijn graphs;
our episode graph (Section~\ref{sec:episodes}) is a distinct object
whose nodes are not residue classes but phantom shadowing episodes.

\subsection{Structure of the paper}
\label{sec:structure}
Section~\ref{sec:setup} fixes notation for the accelerated Syracuse
map and the $\nu_2$-debt formalism. Section~\ref{sec:shadowing}
states and proves the exact congruential shadowing lemma; the
formalization of this section in Lean 4 is summarized in
Section~\ref{sec:formalization}. Section~\ref{sec:episodes}
introduces the episode graph and its critical SCC.
Section~\ref{sec:operator} constructs the finite transfer operator
at the critical node. Section~\ref{sec:bound} develops the weighted
Collatz--Wielandt perturbative bound. Section~\ref{sec:numerics}
reports numerical evidence in $T$ and $j$. Section~\ref{sec:cross}
extends the bound to a cross-node operator on the entire SCC.
Section~\ref{sec:reduction} formulates the explicit conjectures
whose resolution would yield a proof, and includes a quantitative
comparison with Chang \cite{chang2026}. Section~\ref{sec:limits}
is an honest accounting of what this work does and does not
establish. Section~\ref{sec:methodology} documents the AI-assisted
methodology and the contribution split.

\section{Setup: accelerated Syracuse map and gravitational debt}
\label{sec:setup}

For an odd integer $n$, the \emph{Syracuse step} is
\[
S(n) = \frac{3n+1}{2^{\nutwo(3n+1)}},
\qquad
a(n) := \nutwo(3n+1).
\]
We write $a_k = a(n_k)$ for the $\nu_2$ sequence along the orbit
$n_0, n_1, n_2, \ldots$ with $n_{k+1} = S(n_k)$.

A direct expansion gives, after $k$ Syracuse steps,
\[
\log_2 n_k = \log_2 n_0 + k \log_2 3 - \sum_{i=0}^{k-1} a_i + O(1).
\]
The dominant term $k \log_2 3 - \sum a_i$ is what we call the
\emph{gravitational debt} after $k$ steps:
\[
D(k) := k \log_2 3 - \sum_{i=0}^{k-1} a_i.
\]
Positive $D(k)$ indicates net expansion to step $k$; the orbit
height $\log_2 n_k$ tracks $D(k)$ up to bounded fluctuations. The
quantity is otherwise classical and is the standard
\emph{logarithmic drift} or \emph{total stopping-time deficit}
analysis underlying Tao \cite{tao2019} and the historical
literature \cite{lagarias1985,lagarias2010,inselmann2024}; we use
the term ``gravitational debt'' as informal vocabulary for the
program note in Section~\ref{sec:methodology}.

The threshold $a_i > \log_2 3 \approx 1.585$ separates dissipative
from expansive single steps. The Collatz conjecture is morally
equivalent to the assertion that $D(k) \to -\infty$ along every
positive-integer orbit.

\section{Phantom rational orbits and the shadowing lemma}
\label{sec:shadowing}

\subsection{Phantom words and their rational fixed points}
\label{sec:phantom-words}

Fix a finite word
\[
w = (a_0, a_1, \ldots, a_{L-1}) \in \N^L,
\qquad
A = A(w) := \sum_{i=0}^{L-1} a_i.
\]
The composition $S_w := S_{a_{L-1}} \circ \cdots \circ S_{a_0}$ is
affine: writing it out one obtains
\begin{equation}\label{eq:Sw}
S_w(n) = \frac{3^L\, n + C_w}{2^A},
\qquad
\begin{aligned}
C_0 &= 0, \\
C_{j+1} &= 3 C_j + 2^{A_j}, \\
A_{j+1} &= A_j + a_j.
\end{aligned}
\end{equation}
The fixed point of $S_w$ in $\Q$ is
\begin{equation}\label{eq:qw}
q_w = \frac{C_w}{2^A - 3^L}.
\end{equation}
The equation \eqref{eq:qw} is classical: it appears as the
``Bohm--Sontacchi'' criterion in the Collatz literature
\cite[Theorem 1]{siegel2023a} and is a direct consequence of the
parity-vector inverse formula of \cite{lagarias1985,rozier2025}; in
the recent literature it is the equation defining the phantom
cycles of \cite[Def.~7.2]{chang2026}, of which the present
construction is an instance.

\begin{definition}[expansive phantom]
The word $w$ is an \emph{expansive phantom} if $3^L > 2^A$. Then
$2^A - 3^L < 0$, and assuming $C_w > 0$ (which is automatic from
\eqref{eq:Sw}) we have $q_w \in \Q_{<0}$.
\end{definition}

The orbit of $q_w$ under $S$ is exactly periodic of length $L$ with
$\nu_2$ word $w$ repeated. In the $2$-adic completion $\Zp$, $q_w$
makes sense as a $2$-adic integer (since $\gcd(2^A - 3^L, 2)=1$) and
its $2$-adic orbit consists of $L$ distinct $2$-adic integers cycled
by $S$ (lifted to its natural $2$-adic extension; see
Section~\ref{sec:formalization} for the formalized statement).

\subsection{Exact congruential shadowing}

Let $w^\infty$ denote the right-infinite repetition of $w$, and let
$B_m := \sum_{i=0}^{m-1} a_{i \bmod L}$ for $m \ge 0$ be the partial
sum of the periodic word.

\begin{lemma}[Exact shadowing]\label{lem:shadowing}
Let $w$ be an expansive phantom with $2$-adic representative $q_w$,
and let $n$ be a positive odd integer. If
\[
n \equiv q_w \pmod{2^{B_m+1}}
\qquad\text{(equality in $\Zp$ truncated to $B_m+1$ bits),}
\]
then the first $m$ Syracuse steps of $n$ have $\nu_2$ word exactly
$(a_0, a_1, \ldots, a_{m-1})$.

In particular, if $n \equiv q_w \pmod{2^{bA+1}}$, the first $b$
periods (i.e.\ $bL$ Syracuse steps) of $n$ shadow $w^\infty$.
\end{lemma}

\begin{remark}[relation to existing identities]
\label{rem:lemma-context}
Lemma~\ref{lem:shadowing} is the iterated, step-by-step form of
Chang's per-block $2$-adic repulsion identity
$\nu_2(F_\sigma(x) - \rho) = \nu_2(x - \rho) - K$
\cite[Prop.~7.4]{chang2026}, applied $L$ times. Equivalently, it is
the inverse parity-vector formula \cite[Lemma~1]{rozier2025}
restated in terms of the rational fixed point $q_w$ as the
canonical $2$-adic representative of the residue class. The form
chosen here is the one that feeds directly into the episode graph
construction (Section~\ref{sec:episodes}) and into the Lean
formalization (Section~\ref{sec:formalization}).
\end{remark}

\begin{proof}
We argue by induction on $j$. Suppose after $j$ Syracuse steps the
orbits of $n$ and $q_w$ have followed the same word
$(a_0,\ldots,a_{j-1})$. Then both have undergone the same affine
transformation $S_{(a_0,\ldots,a_{j-1})}$, which has multiplier
$3^j / 2^{B_j}$. Therefore
\[
S^j(n) - S^j(q_w) = \frac{3^j}{2^{B_j}} (n - q_w).
\]
Taking $\nu_2$ of both sides and using $\nu_2(3^j)=0$,
\[
\nutwo\bigl(S^j(n) - S^j(q_w)\bigr)
= \nutwo(n - q_w) - B_j
\ge (B_m + 1) - B_j.
\]
For step $j$ we need to extract \emph{exactly} $a_j$ factors of $2$
from $3 S^j(n) + 1$, which agrees with $3 S^j(q_w) + 1$ up to
valuation
\[
\nutwo\bigl(3 S^j(n) + 1 - 3 S^j(q_w) - 1\bigr)
= \nutwo\bigl(3(S^j(n) - S^j(q_w))\bigr)
\ge B_m + 1 - B_j.
\]
For the next valuation to coincide it suffices that
\[
B_m + 1 - B_j \ge a_j + 1,
\quad\text{i.e.}\quad
B_m \ge B_j + a_j = B_{j+1},
\]
which holds for all $j < m$ by definition. Induction completes the
proof. The ``+1'' bit is essential: without it one only ensures
agreement modulo $2^{a_j}$, not the absence of an additional factor
of $2$.
\end{proof}

\begin{corollary}[No infinite shadowing]\label{cor:no-infinite}
A positive integer cannot shadow $w^\infty$ for all $m$, since that
would force $n = q_w$ in $\Zp$, hence $n = q_w \in \Q_{<0}$,
contradicting $n \in \N$. (Compare \cite[Lemma~9.20]{chang2026} for
the per-block form.)
\end{corollary}

\subsection{Formal verification in Lean 4}
\label{sec:formalization}

Lemma~\ref{lem:shadowing} and Corollary~\ref{cor:no-infinite} have
been mechanically verified in the Lean~4 theorem prover with
Mathlib (Lean~v4.29.1, Mathlib~v4.29.1). The formalization is
\texttt{sorry}-free and is included as supplementary material in
the project repository under \texttt{lean/}. The headline
declarations are:
\begin{itemize}
  \item \texttt{CollatzShadowing.exact\_shadowing} ---
        Lemma~\ref{lem:shadowing}, stated paper-faithfully in the
        $2$-adic form: for every positive natural $n$ and every
        \texttt{PhantomWord} $w$, if $n \equiv q_w$
        $\pmod{2^{B_m+1}}$ (encoded as ideal membership in
        $\Zp$, see \texttt{INVENTORY.md} \S 1.4), then the first
        $m$ iterates of $n$ under \texttt{Syracuse2adic}
        produce exactly the periodic word $w$.
  \item \texttt{CollatzShadowing.exact\_shadowing\_periods} ---
        the periodic specialization $m = b \cdot L$, with the
        paper's congruence bound $2^{b A + 1}$.
  \item \texttt{CollatzShadowing.no\_infinite\_period\_congruence\_expansive}
        --- Corollary~\ref{cor:no-infinite} for expansive
        phantoms: no positive natural $n$ can satisfy
        $n \equiv q_w \pmod{2^{b A + 1}}$ for all $b$.
\end{itemize}
The supporting infrastructure includes a paper-faithful
$2$-adic extension of the Syracuse map
$\mathtt{Syracuse2adic}\colon \Zp \to \Zp$, the structural
\texttt{PhantomWord} datatype with positivity constraints, the
discharge of the odd-denominator hypothesis on $q_w$ for every
phantom, the affine difference identity $S^j(n) - S^j(q_w) =
(3^j / 2^{B_j})(n - q_w)$ in $\Q_{(2)}$, the strict-ultrametric
$\nu_2$-stability under proximity, and the matching property of
$q_w$'s own $S$-orbit (so that no auxiliary hypothesis is needed
on the user-supplied $n$ beyond the congruence). The total Lean
codebase is approximately $2{,}300$ lines across seven files
(\texttt{Basic}, \texttt{Phantom}, \texttt{Syracuse2Adic},
\texttt{Auxiliary}, \texttt{Shadowing}, \texttt{NoInfinite},
\texttt{INVENTORY}). Build instructions and a complete inventory
of the Mathlib API used are documented in \texttt{lean/README.md}
and \texttt{lean/CollatzShadowing/INVENTORY.md} respectively.

\paragraph{v3 extensions to the Lean formalization.}
In v3 the formalization is extended beyond the shadowing core to
include the episode graph, the truncated transfer operator with
phase-state quotient, the generic weighted Collatz--Wielandt bound,
two generated finite numerical certificates, and the corresponding
spectral-radius bridges to Mathlib's \texttt{spectralRadius}. The
project remains \texttt{sorry}-free and \texttt{axiom}-free in the
new modules. The relevant new modules and declarations are:
\begin{itemize}
  \item \textbf{Episode graph.}
        \texttt{CollatzShadowing.EpisodeGraph}
        introduces \texttt{EpisodeNode}, the finite cutoff box
        \texttt{TruncatedEpisodeNode K C B}, the edge relation
        \texttt{EpisodeGraph.edge}, reachability and
        strongly-connected-component wrappers, and the concrete
        finite-walk certificate format \texttt{HubSCCCertificate}
        based on paths to and from a hub. The Mathlib API
        inventory is documented in
        \texttt{CollatzShadowing/EPISODE\_INVENTORY.md}.
  \item \textbf{Phase state and operator.}
        \texttt{CollatzShadowing.Operator} defines
        \texttt{PhaseState V := Fin (V+1) $\times$ Fin 4 $\times$
        Fin 4}, the \texttt{cappedNu2}, \texttt{oddPart},
        \texttt{mod4Fin}, \texttt{phaseState} projections, and the
        operator interfaces
        \texttt{TransferMatrix V}, \texttt{ProbabilityEntry},
        \texttt{RowSubstochastic}, and the bundled
        \texttt{OperatorDecomposition} record packaging
        $\full = \core + \tail$ together with the substochastic
        row bounds.
  \item \textbf{Weighted Collatz--Wielandt bound.}
        \texttt{CollatzShadowing.Bound} proves the generic finite
        pointwise bound used in
        Section~\ref{sec:bound}: a
        \texttt{FiniteCWCertificate} packages a positive vector
        $v$, a constant $\alpha$, and a per-row inequality
        $(M v)_i \le \alpha \cdot v_i$, and the bridge
        \ttbr{spectralRadius\_le\_of\_finiteCWCertificate}
        promotes it to
        $\spec(M_\R) \le \alpha$ in Mathlib's real
        \texttt{spectralRadius}. The additive
        $\full = \core + \tail$ form is closed by
        \texttt{FiniteCWCertificate.add},
        \texttt{finiteCWCertificateOfSumEq}, and
        \texttt{OperatorDecomposition.cwCertificate\_full}, with
        the spectral corollary
        \texttt{OperatorDecomposition.spectralRadius\_le\_full}.
  \item \textbf{Generated single-node certificate at $T = 10$,
        $j = 32$.}
        \texttt{CollatzShadowing.Generated.T10CriticalSymbolic}
        imports an exact $T = 10$ critical-symbolic full transfer
        matrix from CSV.
        \texttt{CollatzShadowing.Generated.T10J32HighBitTail}
        imports the exact rational
        majority-signature decomposition
        $\full = \core + \tail$ at $T = 10$, $j = 32$,
        defines the bundled
        \texttt{t10j32HighBitTailDecomposition :
        OperatorDecomposition 13}, and proves the three
        substochastic row certificates by exact
        \texttt{decide}-style arithmetic.
        \texttt{CollatzShadowing.Generated.T10J32HighBitTailCW}
        exposes the single-node spectral-radius bound
        \begin{verbatim}
  t10j32HighBitTailSpectralRadiusBound_97_2000 :
    spectralRadius ℝ
      (nnrealMatrixToReal t10j32HighBitTailMatrix)
    ≤ (((97 : NNReal) / (2000 : NNReal)) : ℝ≥0∞).
        \end{verbatim}
        Equivalently, $\spec(\full_{10,32}) \le 97/2000 = 0.0485$.
        The CW vector is generated from
        \ttbr{high\_bit\_tail\_edges\_T10\_j32.csv} and verified
        in 224 evaluated rows split across the
        \texttt{T10J32HighBitTailCWRows00}$\,$\ldots$\,$\texttt{Rows13}
        modules. There is no \texttt{axiom},
        \texttt{sorry}, or \texttt{admit} in the generated CW
        certificate.
  \item \textbf{Generated deterministic finite-residue-cell
        certificate at $K_0 = 16$.}
        \texttt{CollatzShadowing.Generated.K16S16KSCC} imports the
        generated 37-state compressed $K_0 = 16$ SCC certificate
        (Section~\ref{sec:taxonomy}) as a
        \texttt{HubSCCCertificate} plus
        \texttt{CriticalSCCCertificate}.
        \texttt{CollatzShadowing.Generated.K16S16KExactCWSummary}
        imports the empirical 37-node CW certificate summary used
        as historical baseline (\textbf{empirical} max ratio
        $\approx 0.886$, $\alpha = 89/100$) and produces the full
        generated certificate
        \texttt{k16s16KFiniteCWCertificate}.
        \texttt{CollatzShadowing.Generated.K16S16KBridge}
        checks that the generated SCC and CW certificates share
        the same \texttt{Fin 37} state ordering and labels, and
        exposes
        \texttt{k16s16KSpectralRadiusBound}.
        \texttt{CollatzShadowing.Generated.K16S16KDeterministicCW}
        imports the \emph{deterministic} finite residue-cell
        certificate (Section~\ref{sec:taxonomy}) with the strictly
        improved bound, and exposes the spectral-radius bridge
        \begin{verbatim}
  k16s16KDeterministicGeneratedSpectralRadiusBound :
    spectralRadius ℝ
      (nnrealMatrixToReal k16s16KDeterministicMatrix)
    ≤ (k16s16KDeterministicAlphaNNReal : ℝ≥0∞).
        \end{verbatim}
        Equivalently, $\spec(P_{K,b}) \le 3/4$.
\end{itemize}

\paragraph{Scope of the v3 Lean formalization.}
The Lean formalization now \emph{verifies the shadowing core, a
generic weighted Collatz--Wielandt finite-matrix bound, and two
specific generated finite numerical certificates}; it does not
formalize Conjectures~\ref{conj:T} and~\ref{conj:j}, nor the
conditional reduction theorem
(Theorem~\ref{thm:conditional}), since both involve
infinite suprema and asymptotic decay that are statable but not
presently provable in Mathlib. The bounds carried into
\texttt{spectralRadius} by
\ttbr{t10j32HighBitTailSpectralRadiusBound\_97\_2000} and
\ttbr{k16s16KDeterministicGeneratedSpectralRadiusBound} are
\emph{single-$(T, j)$} and \emph{single-$K_0$} finite statements,
not uniform-in-$T$ statements. We discuss this scope explicitly in
Section~\ref{sec:limits} and in Section~\ref{sec:methodology}.

\section{The episode graph}
\label{sec:episodes}

Lemma~\ref{lem:shadowing} shows that an integer $n$ can shadow a
phantom $w$ for some maximal number of periods $b = b_w(n)$, after
which it must leave the residue class. The natural object for
tracking long orbits of integer rebels is therefore not $n$ itself
but the sequence of \emph{episodes}
\[
n \to (w_1, b_1) \to (w_2, b_2) \to \cdots
\]
together with terminal events (drops below the starting value).

\subsection{The graph}

We monitor a finite set of phantoms identified empirically up to
$k\le 24$ (their explicit list is given in
Appendix~\ref{app:phantoms}). For each pair
$(\text{phantom }k, \text{cycle }c, \text{depth }b)$ we associate a
\emph{node} $(k,c,b)$. The episode graph $G$ has these nodes as
vertices; an edge $(k,c,b) \to (k',c',b')$ records the empirical
observation that some integer orbit shadowed $(k,c,b)$ for exactly
$b$ periods, then re-entered shadowing of $(k',c',b')$.

\begin{remark}[scope vs.\ Phantom Universality]
\label{rem:scope-vs-universality}
The phantom set used here is \emph{empirical}: it is the list of
expansive phantom representatives encountered up to $k \le 24$
during the dense sweeps of Section~\ref{sec:cross}. By the
Phantom Universality theorem of \cite[Theorem~7.15]{chang2026},
\emph{every} primitive cyclic composition of valuations is a
phantom family, so the universe of phantoms is in fact infinite
and grows like $1.87^K$ in $K = \sum k_i$. The empirical sweep
finds no SCC outside the listed nodes for the parameter ranges
explored, but verifying that no further SCC contributes at higher
$K$ is an open question; we discuss this as a missing ingredient
of the conditional reduction in Section~\ref{sec:limits}.
\end{remark}

\subsection{The critical SCC}

Empirical exhaustion (script \verb|59_shadow_episode_graph.py|,
\verb|60_episode_graph_diagnostics.py|) over the dense sample
$b \le 16$, $t \le 65535$ identifies a single nontrivial strongly
connected component containing
\[
(10,1,b) \text{ for } b \in [1,15], \quad
(11,1,b) \text{ for } b \in [1,15], \quad
(12,1,b) \text{ for } b \in [1,2], \quad
(12,2,b) \text{ for } b \in [1,2].
\]
The \emph{critical node} is $(12, 2, 1)$: it is the unique node of
the SCC with significant self-recurrence, and most cross-node
transitions terminate at it.

The strategy of the rest of the paper is to prove that this SCC is
spectrally subcritical when encoded as a transfer operator with
appropriate weights.

\section{The critical transfer operator}
\label{sec:operator}

Fix the critical node $(k,c,b) = (12,2,1)$ with phantom word $w$,
period sum $A$, and $2$-adic representative $q_w$. Every integer $n$
shadowing this node for one period satisfies
\[
n \equiv q_w \pmod{2^{A+1}},
\quad\text{i.e.}\quad
n = r + t \cdot 2^{A+1}
\]
for the unique residue $r = q_w \bmod 2^{A+1}$ and parameter
$t \in \N$.

\subsection{Refined phase state}

For a parameter $t$, we define the \emph{refined phase state}
\[
\sigma(t, h) := \bigl(\nu_2(t)\wedge V,\ \mathrm{odd}(t) \bmod 4,\ h \bmod 4\bigr),
\]
where $V$ is a finite valuation cap, $\mathrm{odd}(t) := t / 2^{\nu_2(t)}$,
and $h$ is the count of monitored hits accumulated so far. The
choice of moduli $4$ and $4$ is the empirical minimum that makes
the spectral truncation stable (script
\verb|74_phase_transfer_stress.py|).

\begin{remark}[difference from Chang's state space]
\label{rem:state-space}
Chang \cite[Sect.~9, Rem.~9.17]{chang2026} works with the transfer
operator $T_K$ on the residue space $(\Z/2^K\Z)_{\mathrm{odd}}$ and
proves it nilpotent: every non-leading eigenvalue vanishes and
iterations absorb distributions onto the periodic orbits of the
map mod $2^K$. The refined phase state above is a \emph{different}
quotient: it retains the $2$-adic valuation of the lift parameter
$t$ (capped at $V$) and a low-resolution view of the odd part of
$t$ together with a hit counter, rather than the full residue mod
$2^K$. As a consequence $\full_{T,j}$ is not nilpotent --- it has
$\spec(\full) \approx 0.03$ at the critical node --- and is
suitable for a finite spectral bound rather than a finite-step
absorption argument.
\end{remark}

\subsection{Truncated finite operator}

For depths $T \in \N$ and $j \in \N$, we define the truncated
transfer operator $\full_{T,j}$ on the state space
\[
\mathcal{S}_{T,j} := \{\sigma(t, h) : t \in [0, 2^T \cdot j),\ h \in [0,4)\}
\]
as follows. For each source state $\sigma_{\mathrm{src}}$, enumerate
all $(t, h)$ producing it. Trace the Syracuse orbit from
$n_0 = r + t \cdot 2^{A+1}$ until either:
\begin{enumerate}
  \item $n < n_0$ (terminal: contributes no edge);
  \item the next monitored hit at the same critical node occurs at
        $n^* = r + t^* \cdot 2^{A+1}$. This contributes an edge
        \[
        \sigma_{\mathrm{src}} \to \sigma(t^*, h+1)
        \quad\text{with weight}\quad
        2^{-\delta},
        \quad
        \delta = \lfloor \log_2 t^* \rfloor - \lfloor \log_2 t \rfloor.
        \]
\end{enumerate}
We aggregate edges and normalize by source counts:
\[
(\full_{T,j})_{\sigma, \sigma'}
= \frac{1}{|\{(t,h) : \sigma(t,h)=\sigma\}|}
  \sum_{\substack{(t,h) \\ \sigma(t,h)=\sigma}}
  \mathbf{1}_{\text{trace lands on }\sigma'}
  \cdot 2^{-\delta(t)}.
\]

By construction $\full_{T,j}$ is a substochastic-like nonnegative
matrix with row sums $\le 1$, the deficit corresponding to terminal
transitions.

\subsection{Empirical signature decomposition}

For each pair $(r, h)$ with $r \in [0, 2^T)$, the $j$ lifts
\[
t = r + j' \cdot 2^T, \quad j' \in [0, j),
\]
produce $j$ traces. We collect the multiset of \emph{transition
signatures} (destination state plus $\delta$ bit increment) and
declare the \emph{core signature} to be the most frequent one. Edges
matching this signature are placed in $\core_{T,j}$; all other
edges are placed in $\tail_{T,j}$. By construction
\[
\full_{T,j} = \core_{T,j} + \tail_{T,j}.
\]
We refer the reader to script
\verb|77_high_bit_tail_bound.py| for the implementation details.
This decomposition --- partitioning by empirical majority of
transition signatures --- is, to our knowledge, not present in
prior work on Collatz transfer operators
\cite{chang2026,neklyudov2022,leventides2020,mori2025}.

\section{Weighted perturbative bound}
\label{sec:bound}

\subsection{Collatz--Wielandt with the Perron eigenvector of FULL}

The matrices $\full_{T,j}$, $\core_{T,j}$, $\tail_{T,j}$ are
nonnegative. Let $v_{T,j}$ denote a right Perron eigenvector of
$\full_{T,j}$, normalized so $\|v_{T,j}\|_\infty = 1$. By
Perron--Frobenius (after a tiny rank-1 perturbation to ensure
irreducibility, see remark below), $v_{T,j}$ exists and is strictly
positive.

The decomposition $\full_{T,j} = \core_{T,j} + \tail_{T,j}$ acting
on $v_{T,j}$ gives, componentwise,
\[
(\full_{T,j} v_{T,j})_i
= (\core_{T,j} v_{T,j})_i + (\tail_{T,j} v_{T,j})_i.
\]
Taking the maximum over $i$ such that $(v_{T,j})_i > 0$ and
applying the Collatz--Wielandt characterization of the spectral
radius for non-negative matrices \cite{lemmensnussbaum,
akiangaubert} gives
\begin{equation}\label{eq:bound}
\boxed{\;
\spec(\full_{T,j})
\;\le\;
\underbrace{\max_i \frac{(\core_{T,j} v_{T,j})_i}{(v_{T,j})_i}}_{\mathrm{core\ action}}
\;+\;
\underbrace{\max_i \frac{(\tail_{T,j} v_{T,j})_i}{(v_{T,j})_i}}_{\mathrm{tail\ action}}.
\;}
\end{equation}
We denote the right-hand side $\mathrm{bound}(T,j)$.

\begin{remark}[numerical implementation]
We compute $v_{T,j}$ by power iteration on
$\full_{T,j} + \varepsilon\, J$ with $\varepsilon = 10^{-12}$ and
$J$ the all-ones matrix. The rank-1 perturbation removes
reducibility numerically without affecting the dominant eigenvalue
beyond $\varepsilon$. Convergence is checked in $\ell^\infty$ norm
to tolerance $10^{-12}$.
\end{remark}

\begin{remark}[why CORE's Perron is wrong]
Using the Perron eigenvector of $\core_{T,j}$ instead of
$\full_{T,j}$ yields a divergent ``bound'': $\core_{T,j}$ is
typically reducible, with its Perron eigenvector concentrated on a
strict subset of states, and $\tail_{T,j}$ routinely sends mass to
states outside that subset. The Collatz--Wielandt ratios then blow
up. The empirical fix and the theoretical fix both point to using
$v_{T,j}$ from $\full_{T,j}$.
\end{remark}

\section{Numerical evidence}
\label{sec:numerics}

\subsection{Bound stability in $T$}

We compute $\mathrm{bound}(T, j)$ for $T \in \{8, 10, 12, 14, 16\}$
and $j \in \{16, 32\}$ using script
\verb|84_weighted_perturbation_bound.py| (with the eigenvector
selection from the previous remark). Results:

\begin{center}
\begin{tabular}{rrrrrr}
\toprule
$T$ & $j$ & $\spec(\full)$ & $\mathrm{core\ action}$ & $\mathrm{tail\ action}$ & $\mathrm{bound}$ \\
\midrule
 8 & 16 & 0.0294 & 0.0114 & 0.0294 & 0.0408 \\
 8 & 32 & 0.0299 & 0.0102 & 0.0299 & 0.0400 \\
10 & 16 & 0.0303 & 0.0184 & 0.0303 & 0.0487 \\
10 & 32 & 0.0308 & 0.0177 & 0.0308 & 0.0485 \\
12 & 16 & 0.0303 & 0.0137 & 0.0303 & 0.0499 \\
12 & 32 & 0.0306 & 0.0135 & 0.0306 & 0.0497 \\
14 & 16 & 0.0306 & 0.0157 & 0.0306 & 0.0511 \\
14 & 32 & 0.0307 & 0.0154 & 0.0307 & 0.0507 \\
16 & 16 & 0.0307 & 0.0169 & 0.0307 & 0.0516 \\
\bottomrule
\end{tabular}
\end{center}

The bound is monotonically increasing in $T$ but the increments
decay geometrically:
\[
\Delta_8^{10} = +0.0079, \quad
\Delta_{10}^{12} = +0.0012, \quad
\Delta_{12}^{14} = +0.0012, \quad
\Delta_{14}^{16} = +0.0005.
\]
A geometric extrapolation yields
\[
\lim_{T \to \infty} \mathrm{bound}(T, 16) \le 0.054.
\]
The bound is essentially independent of $j$: differences between
$j=16$ and $j=32$ at fixed $T$ never exceed $0.001$.

\subsection{Spectral gap $\spec(\full) - \spec(\core)$}

Independently, the spectral radii of $\full$ and $\core$ are
computed (script \verb|83_perturbation_gap_scaling.py|). The gap
$\spec(\full) - \spec(\core)$ shrinks with $T$:
\[
T=6:\ 0.025\text{--}0.032,
\quad
T=8:\ 0.025\text{--}0.028,
\quad
T=10:\ 0.018\text{--}0.019,
\]
\[
T=12:\ 0.016\text{--}0.017,
\quad
T=14:\ 0.015,
\quad
T=16:\ 0.014.
\]
The spectral mass of $\full$ migrates progressively to $\core$ as
the quotient is refined.

\subsection{Truncation stability in $j$ (Cauchy property)}

For fixed $T$, we check whether $\full_{T, j}$ admits a Cauchy
property in $j$ via two independent metrics
(script \verb|85_truncation_stability.py|):
\begin{itemize}
  \item \textbf{signature drift}: the fraction of $(r, h)$ groups
        whose dominant signature changes between consecutive
        doublings of $j$;
  \item \textbf{bound increment}: $\mathrm{bound}(T, 2j) - \mathrm{bound}(T, j)$.
\end{itemize}

Results:

\begin{center}
\begin{tabular}{rrrrrr}
\toprule
$T$ & $j$ & sig\_drift\_vs\_prev & bound & $\Delta_{\mathrm{prev}}$ \\
\midrule
 8 & 16  & --     & 0.0408 & --      \\
 8 & 32  & 0.0078 & 0.0400 & $-0.0007$ \\
 8 & 64  & 0.0078 & 0.0421 & $+0.0021$ \\
 8 & 128 & 0.0000 & 0.0425 & $+0.0004$ \\
\midrule
10 & 16  & --     & 0.0487 & --      \\
10 & 32  & 0.0039 & 0.0485 & $-0.0003$ \\
10 & 64  & 0.0088 & 0.0509 & $+0.0024$ \\
10 & 128 & 0.0000 & 0.0515 & $+0.0006$ \\
\bottomrule
\end{tabular}
\end{center}

Two facts:
\begin{enumerate}
  \item The signature drift drops to \emph{exactly zero} when $j$ is
        doubled from $64$ to $128$, for both $T \in \{8, 10\}$. The
        empirical transition signatures of all $(r, h)$ groups
        stabilize.
  \item The bound increments decay by a factor $\sim 4$ per
        doubling, suggesting a Cauchy bound
        $\lim_{j\to\infty}\mathrm{bound}(T=10, j) \lesssim 0.052$.
\end{enumerate}
Combining the two extrapolations, we conjecture
\[
\sup_{T,\, j} \mathrm{bound}(T, j) < 0.06,
\]
a margin of $16\times$ from the criticality threshold $1$.

\section{Cross-node certificate on the SCC}
\label{sec:cross}

The single-node bound at the critical node is the worst case:
running script \verb|86_scc_node_certificates.py| at $T=10, j=32$
across the SCC nodes yields:

\begin{center}
\begin{tabular}{lrrr}
\toprule
node & $\spec(\full)$ & $\mathrm{bound}$ & status \\
\midrule
$(10,1,1)$ & $\approx 0$ & $0.0000$ & OK (return-to-self vacuous) \\
$(10,1,2)$ & $\approx 0$ & $0.0000$ & OK \\
$(11,1,1)$ & $\approx 0$ & $0.0000$ & OK \\
$(11,1,2)$ & $\approx 0$ & $0.0000$ & OK \\
$(12,1,1)$ & $0.0035$  & $0.0043$ & OK \\
$(12,2,1)$ & $0.0308$  & $0.0485$ & OK (worst case) \\
$(12,2,2)$ & $0.0007$  & $0.0007$ & OK \\
\bottomrule
\end{tabular}
\end{center}

The nodes $(10, 1, *)$, $(11, 1, *)$, $(12, 2, 2)$ have nearly
vanishing return-to-self spectrum. They are entry gates to the SCC,
not self-recurrent.

\subsection{The cross-node operator}

We construct the assembled cross-node operator $M_{\mathrm{cross}}$
on the union state space
\[
\bigsqcup_{(k,c,b) \in \mathrm{SCC}} \mathcal{S}_{T,j}^{(k,c,b)}.
\]
Edges are recorded between the source state at one node and the
\emph{first hit} state at any node of the SCC, with weight
$2^{-\delta_{\mathrm{bit}}}$ where $\delta_{\mathrm{bit}}$ is the
bit-length increment between source $n_0$ and destination value (a
node-invariant quantity). The construction is implemented in
script \verb|87_cross_node_transfer.py|.

At $T=8, j=8$ on $5$ representative SCC nodes:
\[
\spec(M_{\mathrm{cross}}) = 0.0366, \qquad
\mathrm{bound}(M_{\mathrm{cross}}) = 0.0366.
\]
The cross-coupling \emph{lowers} the spectrum from the single-node
worst case ($0.0485$). Topologically the SCC is a star centered on
$(12,2,1)$:
\begin{verbatim}
    (10,1,1) -> (12,2,1)   8128  (gate)
    (11,1,1) -> (12,2,1)   8192  (gate)
    (12,2,2) -> (12,2,1)   8128  (gate)
    (12,2,1) -> (12,2,1)    812  (only significant self-loop)
    (12,2,1) -> (12,1,1)    308
    (12,1,1) -> (12,1,1)     96
    (12,2,2) -> (12,2,2)     64
    (12,2,1) -> (12,2,2)      8
\end{verbatim}
The single-node bound was therefore overly conservative; including
cross-node transitions strictly improves it.

\paragraph{Empirical scope of this section.}
The five-phantom set $\{(10,1,*), (11,1,*), (12,1,*), (12,2,*), (20,1,*)\}$
used in this section is the empirical set encountered during the
dense sweeps of v1 and v2 (Appendix~\ref{app:phantoms}). The
\emph{universe} of expansive phantoms is in fact infinite by
\cite[Theorem~7.15]{chang2026}, growing like $1.87^K$ in
$K = \sum k_i$. Verifying that no further phantom family at higher
depths contributes a separate SCC was left as an open scope issue
in v2. In v3 we close that scope at the explicit depth cutoff
$K_0 = 16$, derive a deterministic finite-residue-cell
Collatz--Wielandt certificate on the resulting taxonomy SCC, and
import it into the Lean formalization. The construction is the
subject of Section~\ref{sec:taxonomy}.

\section{Phantom-set taxonomy and the deterministic
finite-residue-cell certificate at $K_0 = 16$}
\label{sec:taxonomy}

This section extends the empirical SCC of Section~\ref{sec:cross} to
the full primitive expansive phantom universe up to depth
$K_0 = 16$, identifies the resulting taxonomy SCC, certifies a
deterministic finite-residue-cell Collatz--Wielandt bound on the
compressed macro-state space $(K, b)$, and packages the certificate
as a Lean theorem under
\texttt{CollatzShadowing.Generated.K16S16KDeterministicCW}.

\subsection{Necklace enumeration of expansive compositions}
\label{sec:necklace}

By M\"obius inversion of the necklace counts
\cite[\S 2]{chang2026}, the number of primitive cyclic compositions
$(k_1, \ldots, k_\ell)$ of $K$ into $\ell$ positive parts is
\[
M(K, \ell)
= \frac{1}{\ell}
  \sum_{d \mid \gcd(K, \ell)} \mu(d)
  \binom{K/d - 1}{\ell/d - 1}.
\]
Script
\verb|scripts/phantom_taxonomy/necklace_counts.py|
implements $M(K, \ell)$, brute-force verifies the formula for
$K \le 10$, and reproduces the displayed $R(K)$ values of
\cite[\S 7.3]{chang2026} for $K = 3, \ldots, 20$
(supplementary file
\verb|necklace_counts_k3_20.csv|).

\subsection{Phantom representatives in exact arithmetic}
\label{sec:representatives}

For each primitive expansive composition
$w = (k_1, \ldots, k_\ell)$ with drift
$\Delta = \ell \log_2 3 - K > 0$, the rational fixed point
$q_w = C_w / (2^K - 3^\ell)$ is computed in exact arithmetic by
script
\verb|scripts/phantom_taxonomy/phantom_representatives.py|. The
output table
\verb|phantom_representatives_k3_16.csv|
records, for each expansive composition,
$\bigl(w, C_w, K, \ell, q_w, q_w \bmod 2^m\bigr)$.
Three sanity checks are enforced per row: the equation
$(2^K - 3^\ell)\, q_w = C_w$, the fixed-point property
$S_w(q_w) = q_w$, and odd-denominator of $q_w$ in lowest terms
(consistent with the discharge of \texttt{QwOddDen} in
\texttt{Auxiliary.lean}). The per-$K$ row counts match the
M\"obius enumeration $M_{\mathrm{expanding}}(K)$ from
Section~\ref{sec:necklace}. The total number of expansive
representatives with $K \le 16$ is $1247$.

\subsection{Orbit harness and the $K_0 = 16$ taxonomy SCC}
\label{sec:harness}

For each phantom representative, the orbit harness
\verb|scripts/phantom_taxonomy/orbit_harness.py| samples integer
lifts $n = q_w \bmod 2^{bA + 1} + t \cdot 2^{bA + 1}$, traces odd
Syracuse orbits, monitors phantom-congruence hits along the way,
and emits per-source CSVs of nodes, events, and edges. Running
this at $K_0 = 16$ with $b \in \{1, 2\}$, eight dense lifts per
source, and a maximum of $1000$ Syracuse steps per orbit
(supplementary output prefix
\verb|orbit_harness_k16_full_*|) yields:
\begin{itemize}
  \item $2401$ observed taxonomy nodes;
  \item $5041$ distinct edge types in the event stream;
  \item a single nontrivial strongly connected component of size
        $1222$ raw states (full node list in
        \verb|orbit_harness_k16_full_scc_nodes.csv|, summary in
        \verb|notes/phantom_taxonomy_k16_scc_report.md|).
\end{itemize}

\subsection{Compression to $(K, b)$ and the empirical baseline}
\label{sec:empirical-baseline}

To stabilize the Collatz--Wielandt analysis, we compress the raw
1240-node observed source space along two macro projections,
$(K, L, b)$ and $(K, b)$, and aggregate edges accordingly. The
substochastic retention matrix on each compression has been
empirically certified via
\verb|scc_transfer_summary.py| /
\verb|scc_cw_certificate.py| /
\verb|scc_collatz_wielandt.py|. With $16$ lifts per source the
empirical maximum Collatz--Wielandt ratios are:
\begin{center}
\begin{tabular}{lrr}
\toprule
compression & nodes & exact CW max ratio \\
\midrule
raw node      & $1240$ & $439764459109/496636575879   \approx 0.88548$ \\
$(K,L,b)$     & $76$   & $88036787882257/99446226949575 \approx 0.88527$ \\
$(K,b)$       & $37$   & $136756256754601/154382162832639 \approx 0.88583$ \\
\bottomrule
\end{tabular}
\end{center}
This certifies sampled empirical matrices and is retained in v3 as
the \emph{historical baseline}. The Lean import of this empirical
$(K, b)$ certificate is
\texttt{CollatzShadowing.Generated.K16S16KExactCWSummary}, with
$\alpha = 89/100$
(\texttt{k16s16KFiniteCWCertificate}).

\subsection{Deterministic finite residue-cell certificate}
\label{sec:deterministic}

The empirical certificate of Section~\ref{sec:empirical-baseline}
certifies sampled transition matrices. We strengthen this to a
\emph{deterministic} construction at the same depth $K_0 = 16$.
For each of the $1240$ raw SCC source states, the deterministic
generator
\verb|scripts/phantom_taxonomy/deterministic_residue_transfer.py|
enumerates all $2^{\mathrm{lift\_bits}}$ finite residue subclasses
of the lift parameter $t$ (the small residue refinement above
$2^{A + 1}$) and classifies each cell exactly:
\begin{itemize}
  \item \emph{canonical source cell}: a cell that initiates a
        deterministic transition to a definite destination cell
        under one shadowing period;
  \item \emph{shadowed initial cell}: a cell that descends
        immediately to a covered destination of a canonical
        source;
  \item \emph{exit below start}: a cell whose deterministic
        trace falls below $n_0$ within the budget;
  \item \emph{budget exit}: a cell whose deterministic trace
        exceeds the step budget. With
        $\mathrm{lift\_bits} = 4$ and step budget $1000$, the
        number of budget exits is \textbf{zero}.
\end{itemize}
The enumeration produces exact rational transition matrices on
each of the macro projections (raw node, $(K,L,b)$, $(K,b)$).
Coverage figures at $\mathrm{lift\_bits} = 4$:
\begin{center}
\begin{tabular}{rrrrr}
\toprule
total cells & canonical source & shadowed initial &
no-initial & budget exits \\
\midrule
$19840$ & $17671$ & $2169$ & $0$ & $0$ \\
\bottomrule
\end{tabular}
\end{center}
The retained certified macro-state space is $(K, b)$:
\begin{center}
\begin{tabular}{rrrrr}
\toprule
states & nonzero internal edge types &
source events & internal transitions & exits below start \\
\midrule
$37$ & $182$ & $17671$ & $17176$ & $495$ \\
\bottomrule
\end{tabular}
\end{center}

The exact Collatz--Wielandt certificate is verified by
\verb|scc_cw_certificate.py --verify|:
\[
\boxed{\;
  \max_i \frac{(P_{K,b}\, v)_i}{v_i}
  = \frac{90833233962213}{129559208330288}
  \approx 0.70109
  \;<\;
  \frac{3}{4},
  \;}
\]
with max-ratio node $K11\!:\!b2$. The CW witness vector $v$ is a
positive integer vector; the inequality holds rowwise in exact
rational arithmetic.

\paragraph{Sensitivity at finer residue refinement.}
The same generator at
$\mathrm{lift\_bits} \in \{4, 5, 6\}$ produces:
\begin{center}
\begin{tabular}{rrrrl}
\toprule
$\mathrm{lift\_bits}$ & source cells & canonical & exits &
exact max ratio \\
\midrule
$4$ & $19840$ & $17671$ & $495$  &
$90833233962213/129559208330288 \approx 0.70109$ \\
$5$ & $39680$ & $35331$ & $990$  &
$7332495524923/10616480126384 \approx 0.69067$ \\
$6$ & $79360$ & $70667$ & $1975$ &
$64869145309473/97226913303232 \approx 0.66719$ \\
\bottomrule
\end{tabular}
\end{center}
All three remain strictly below $3/4$, and the bound is
monotonically decreasing under finer residue refinement. No
budget exits at any of the three settings.

\subsection{Lean import of the deterministic certificate}
\label{sec:deterministic-lean}

The deterministic certificate is imported into the Lean
formalization via the script
\verb|scripts/phantom_taxonomy/lean_cw_summary.py|, generating
\texttt{CollatzShadowing.Generated.K16S16KDeterministicCW.lean}.
The headline generated declarations are:
\begin{itemize}
  \item
  \texttt{k16s16KDeterministicMatrix} ---
  the $37 \times 37$ non-negative
  rational transition matrix on the
  $(K, b)$ macro-state space, packaged as
  \texttt{Matrix K16S16KDeterministicState
  K16S16KDeterministicState NNReal};
  \item
  \texttt{k16s16KDeterministicCWBasis} --- the positive integer
  CW witness vector (\texttt{FiniteCWBasis});
  \item
  \ttbr{k16s16KDeterministicAlphaNNReal} --- the bound
  $\alpha = 3/4$;
  \item
  \ttbr{k16s16KDeterministicMaxRatio\_lt\_alpha} ---
  the exact rational inequality
  $90833233962213/129559208330288 < 3/4$ proved by
  \texttt{norm\_num};
  \item
  \ttbr{k16s16KDeterministicFiniteCWCertificate} ---
  the full generated finite Collatz--Wielandt certificate
  (\texttt{FiniteCWCertificate});
  \item
  \ttbr{k16s16KDeterministicGeneratedSpectralRadiusBound} ---
  the spectral-radius bridge to Mathlib's real
  \texttt{spectralRadius}:
  \begin{verbatim}
  spectralRadius ℝ
    (nnrealMatrixToReal
      k16s16KDeterministicMatrix)
    ≤ (k16s16KDeterministicAlphaNNReal : ℝ≥0∞).
  \end{verbatim}
\end{itemize}
The generated module is checked under
\verb|lake build CollatzShadowing.Generated.K16S16KDeterministicCW|
with no \texttt{axiom}, \texttt{sorry}, or \texttt{admit}.

\subsection{Scope of the taxonomy certificate}
\label{sec:taxonomy-scope}

The certificate of Section~\ref{sec:deterministic} is the
deterministic finite-residue-cell Collatz--Wielandt bound
\[
  \spec(P_{K, b}^{(K_0 = 16,\, \mathrm{lift\_bits})}) \le 3/4
\]
for $\mathrm{lift\_bits} \in \{4, 5, 6\}$. Three caveats apply.
First, the construction is at the declared depth cutoff
$K_0 = 16$: a separate taxonomy run at $K_0 = 20$ has not been
carried out in v3 and is not required by the present scope (the
$K_0 = 16$ sweep already exhibits the single dominant nontrivial
SCC; see Section~\ref{sec:harness}). Second, the construction
operates on the compressed $(K, b)$ macro projection; the raw
1240-node matrix and the intermediate $(K, L, b)$ projection are
retained as diagnostic certificates with comparable bounds
(Section~\ref{sec:empirical-baseline}). Third, the bound is on
the finite substochastic retention matrix; it is \emph{not} a
spectral-gap proof on an infinite operator. The relationship to
the open analytic question is the subject of
Section~\ref{sec:reduction} and
Section~\ref{sec:limits}.

\section{Conjectures and the spectral reduction statement}
\label{sec:reduction}

We isolate the open analytic content as two explicit conjectures.

\begin{conjecture}[uniform bound in $T$]\label{conj:T}
There exists a constant $B_\infty < 1$ such that
\[
\sup_{T \in \N,\ j \in \N} \mathrm{bound}(T, j) \le B_\infty.
\]
Empirically the data support $B_\infty \approx 0.06$.
\end{conjecture}

\begin{conjecture}[signature closure in $j$]\label{conj:j}
For each $(T, r, h)$ there exists $j^*(T, r, h)$ such that the
dominant transition signature stabilizes for all $j \ge j^*$, and
furthermore
\[
\sup_{j \ge j^*}
\bigl|\mathrm{bound}(T, j) - \mathrm{bound}(T, j^*)\bigr|
\le C \cdot 2^{-\beta(j - j^*)}
\]
for some $C, \beta > 0$ uniform in $(T, r, h)$.
\end{conjecture}

\begin{theorem}[conditional reduction, sketch]\label{thm:conditional}
Assume Conjectures~\ref{conj:T} and~\ref{conj:j}. Then for every
positive integer $n$, the Syracuse orbit $S^k(n)$ satisfies
$\liminf_{k\to\infty} S^k(n) = 1$.
\end{theorem}

The argument is the standard one for transfer-operator subcriticality
combined with the no-infinite-shadowing
Corollary~\ref{cor:no-infinite}: assuming the conjectures, the
critical SCC has spectral radius $\le B_\infty < 1$ on a refined
quotient that captures all dynamically relevant phases, hence the
return-to-critical mass strictly contracts, hence orbits cannot
remain forever near the SCC, hence by exhaustion of the empirical
phantom set every orbit eventually escapes to a dissipative regime
with average $\nu_2 > \log_2 3$ and descends. The full argument
requires an exhaustive enumeration of phantoms below a controlled
threshold and a verification that no further SCCs appear; see
Section~\ref{sec:limits} for an honest discussion of the gaps.

\paragraph{What the conditional reduction would buy.}
The conclusion of Theorem~\ref{thm:conditional},
$\liminf_{k\to\infty} S^k(n) = 1$ for \emph{every} positive
integer $n$, is the full Collatz convergence statement (one
direction of the conjecture). This is qualitatively stronger than
the unconditional ``almost-all'' result of Tao
\cite{tao2019} --- which asserts only logarithmic-density
convergence to almost bounded values. The price of the stronger
conclusion is precisely the open spectral content of
Conjectures~\ref{conj:T} and~\ref{conj:j}: a uniform finite-matrix
bound and a Cauchy-type signature closure on the critical SCC.
Closing the conjectures unconditionally would therefore push the
problem from ``almost-all density'' to ``every $n$'', without
addressing the trivial-cycle uniqueness question separately
treated in the cycle-exclusion literature
(e.g.\ \cite{chang2026}, sect.~v7 cycle filter).

\subsection{Comparison with concurrent work}
\label{sec:comparison}

Chang \cite{chang2026} obtains, by an entirely different route, two
quantitative bounds on the per-orbit phantom contribution that are
conceptually comparable to the spectral bounds developed here:
\begin{itemize}
  \item the unconditional spectral bound on the depth-$2$ extended
        return kernel \cite[Theorem~C.3]{chang2026},
        $\spec(\widetilde{B}_2^{\mathrm{ext}}) \le 5/32 \approx 0.156$,
        with associated per-return information cost
        $\ge \log_2(32/5) \approx 2.678$ bits;
  \item the per-orbit phantom gain rate
        \cite[Theorem~7.19]{chang2026},
        $R \le 0.0893 < \varepsilon = 2 - \log_2 3 \approx 0.415$,
        with safety margin $\varepsilon / R \ge 4.65\times$.
\end{itemize}
The bounds and ours are computed on different objects, so
comparison should be qualitative. Chang's $\widetilde{B}_2^{\mathrm{ext}}$
acts on $(\Z/64\Z)_{\mathrm{odd}}$ projected onto the
five-element invariant core $I_2 = \{7, 27, 31, 59, 63\} \pmod{64}$,
and $R$ aggregates the expanding-drift contribution of every
primitive cyclic composition (an exponentially growing phantom
universe). Our $\full_{T,j}$ acts on a finer phase quotient at a
single critical node and yields
$\spec(\full) \approx 0.0308$ at $T = 16, j = 32$ with weighted
upper bound $\mathrm{bound}(T,j) \le 0.0516$. The cross-node
operator $M_{\mathrm{cross}}$ further sharpens the SCC-level
estimate to $0.0366$.

\begin{center}
\begin{tabular}{lcl}
\toprule
quantity & value & object \\
\midrule
$\spec(\full)$, this work, $T=16$, $j=32$ & $0.0308$ & critical-node single-node \\
$\mathrm{bound}(T,j)$, this work & $\le 0.0516$ & weighted Collatz--Wielandt \\
$\spec(\full_{10,32})$, this work, Lean v3 & $\le 97/2000 = 0.0485$ &
critical-node, Lean-checked \\
$\spec(M_{\mathrm{cross}})$, this work & $0.0366$ & cross-node on SCC \\
$\spec(P_{K,b})$, this work v3 & $\le 3/4$ &
$K_0 = 16$ deterministic finite residue-cell \\
$R$, Chang \cite[Thm.~7.19]{chang2026} & $\le 0.0893$ & per-orbit phantom gain \\
$\spec(\widetilde{B}_2^{\mathrm{ext}})$, Chang \cite[Thm.~C.3]{chang2026} & $\le 5/32 \approx 0.156$ & depth-$2$ I$_2$ kernel \\
\bottomrule
\end{tabular}
\end{center}

\paragraph{v3 update.}
The two new finite Lean-checked bounds added in v3 are bounds on
\emph{different objects} from the empirical numerical bound on
$\full_{T,j}$:
\begin{itemize}
  \item $\spec(\full_{10,32}) \le 97/2000$ is the formally verified
        single-$(T, j)$ bound on the empirical signature
        decomposition of $\full$ at $T = 10$, $j = 32$. It is the
        same numerical object as the $0.0485$ entry in
        Section~\ref{sec:numerics}, now expressed as an exact
        rational and exposed to Mathlib's \texttt{spectralRadius}.
  \item $\spec(P_{K, b}) \le 3/4$ is a bound on the deterministic
        finite-residue-cell substochastic retention matrix on the
        $K_0 = 16$ taxonomy SCC, on the compressed $(K, b)$
        macro-state space. It is \emph{coarser} than the single-node
        $\full$ bound because the taxonomy SCC absorbs many phantom
        families at once, but it is \emph{strictly subcritical} and
        \emph{deterministic} (i.e.\ not based on sampling).
\end{itemize}
The two bounds are complementary: $97/2000$ certifies that the
critical single-node operator is far from the unit circle at
$T = 10$, $j = 32$; $3/4$ certifies that the deterministic enriched
$K_0 = 16$ taxonomy retention is subcritical even after enlarging
the phantom universe far beyond the v2 empirical set. Neither is a
uniform-in-$T$ statement and neither closes
Conjecture~\ref{conj:T}.

The two programs differ in coverage and granularity. Chang's
bounds are unconditional and aggregate but coarser; ours are
finer and local but conditional on the conjectures of this
section. The fact that two independently developed AI-assisted
programs arrive, by very different routes, at quantitatively
similar bounds is, in our view, evidence that the underlying
phenomenology is robust.

\paragraph{Scale of parameters explored.}
The two analyses operate at very different parameter scales, and
the comparison should be read with this in mind. Chang's exact
finite computation of the per-orbit gain rate
\cite[Theorem~7.19, Step 2]{chang2026} sums $R(K)$ for
$K = 3, \ldots, 500$ via M\"obius inversion of the primitive
cyclic-composition counts $M(K, \ell)$, with an analytic
Chernoff--Cram\'er tail beyond $K_0 = 500$, aggregating over a
phantom universe whose size grows like $1.87^K$. The numerical
evidence developed here covers $T \le 16$ and $j \le 128$ at the
single critical node $(12, 2, 1)$, with full phase-state resolution
$(\nu_2(t), \mathrm{odd}(t) \bmod 4, h \bmod 4)$. The two regimes
are complementary --- broad-shallow aggregation over an
exponentially growing family count versus narrow-deep resolution
at a single node --- and not directly substitutable. Extending
the present sweep to higher $T$ is computationally feasible (the
matrix sizes scale polynomially in $T$ and linearly in $j$) and
is part of the planned program for a future revision; we note,
however, that the geometric decay of bound increments observed
through $T = 16$ already suggests a finite limit, so the principal
analytic content lies in proving that limit, not in extending the
sweep numerically.

\paragraph{The same barrier in two formulations.}
Conjectures~\ref{conj:T} and~\ref{conj:j} are a finite-matrix
re-expression of the same distributional-to-pointwise gap that
\cite[Sect.~9]{chang2026} identifies as the irreducible obstruction
to closing every standard route. Concretely, Chang's gain-observable
recurrence \cite[Conj.~9.18]{chang2026},
\[
\eta_{K+3}(n_0) \le \lambda \eta_K(n_0) + C K^{-\beta},
\]
is a depthwise sharpening of the same kind of bound that
Conjecture~\ref{conj:T} requires uniformly in $T$. We do not claim
that closing one closes the other --- the operators and observables
are different --- but the structural barrier is the same. We
therefore adopt the position of \cite{chang2026}: \emph{the present
work is not a proof of the Collatz conjecture; it characterizes
why a particular structural reduction does not close
unconditionally}.

\section{What this is and what this is not}
\label{sec:limits}

This work is a \emph{program}, not a proof. We summarize what is
established, what is empirical, and what is missing.

\paragraph{Established (rigorously, with mechanical verification).}
\begin{itemize}
  \item Lemma~\ref{lem:shadowing} (exact congruential shadowing)
        is a theorem with elementary proof. It identifies the
        precise residue condition for $b$-fold phantom shadowing.
        \emph{It has been mechanically verified in Lean~4 with
        Mathlib} (Section~\ref{sec:formalization}).
  \item Corollary~\ref{cor:no-infinite} (no infinite shadowing)
        follows formally and has also been verified in Lean~4 in
        the form \texttt{no\_infinite\_period\_congruence\_expansive}.
  \item The construction of the finite transfer operators
        $\full_{T,j}$, $\core_{T,j}$, $\tail_{T,j}$ at the critical
        node is a deterministic algorithm; the matrices are
        exactly computable for all finite $T, j$.
  \item The bound \eqref{eq:bound} is a Collatz--Wielandt
        consequence of nonnegativity \cite{lemmensnussbaum}; it
        gives a true upper bound on $\spec(\full_{T,j})$ for each
        finite $(T, j)$.
\end{itemize}

\paragraph{Empirical (numerically supported, not proved).}
\begin{itemize}
  \item All numerical bounds reported in
        Sections~\ref{sec:numerics} and~\ref{sec:cross}.
  \item The geometric decay of bound increments in $T$ and $j$.
  \item The closure of dominant signatures at $j = 128$ for
        $T \in \{8, 10\}$.
  \item The identification of $(12, 2, 1)$ as the worst-case node
        of the SCC.
  \item The structural ``star'' topology of the SCC around
        $(12, 2, 1)$.
\end{itemize}

\paragraph{Closed in v3 relative to v2.}
\begin{itemize}
  \item The taxonomic mapping of the empirical v2 phantom set onto
        the universe of primitive cyclic compositions is carried
        out at the explicit cutoff $K_0 = 16$ in
        Section~\ref{sec:taxonomy}: all $1247$ primitive expansive
        compositions with $K \le 16$ are enumerated by M\"obius
        inversion of necklace counts, their $2$-adic
        representatives are computed in exact arithmetic, and the
        orbit harness produces a single nontrivial taxonomy SCC of
        $1222$ raw states.
  \item The empirical scope of Section~\ref{sec:cross} is
        strengthened by the deterministic finite-residue-cell
        Collatz--Wielandt certificate on the compressed $(K, b)$
        macro-state space (Section~\ref{sec:deterministic}): exact
        rational max ratio $90833233962213/129559208330288 < 3/4$,
        with sensitivity at $\mathrm{lift\_bits} \in \{4, 5, 6\}$
        confirming the bound is monotonically improving under
        finer residue refinement.
  \item The episode graph, the truncated transfer operator with
        phase-state quotient, the empirical signature
        decomposition $\full = \core + \tail$, and the generic
        weighted Collatz--Wielandt bound are formalized in Lean
        (Section~\ref{sec:formalization}), and two specific
        generated finite numerical certificates are imported and
        their spectral-radius bridges are exposed.
\end{itemize}

\paragraph{Still missing.}
\begin{itemize}
  \item Proofs of Conjectures~\ref{conj:T} and~\ref{conj:j}. We
        believe the natural framework is transfer-operator analysis
        on a suitable Banach space of $2$-adic Lipschitz functions;
        a Lasota--Yorke inequality for the family
        $\{\full_T\}$, combined with Hennion's theorem and
        Keller--Liverani perturbation theory \cite{kellerliverani},
        would yield Conjecture~\ref{conj:T}. We have not carried
        this out, and recent literature on quasi-compactness
        \cite{quasicompact} and certified spectral approximation
        \cite{certifiedspectral} suggests this is a non-trivial
        analytic project. The shared barrier with
        \cite[Sect.~9]{chang2026} (the distributional-to-pointwise
        gap) suggests that closing it conditionally on a single
        operator may itself be hard. The deterministic
        finite-residue-cell certificate of
        Section~\ref{sec:deterministic} addresses the
        \emph{taxonomy completeness} side of the original v2 gap
        but does \emph{not} address the asymptotic-in-$T$ side.
  \item A formal version of Theorem~\ref{thm:conditional}. The
        finite-matrix layer of the construction is now formalized
        (Section~\ref{sec:formalization}); the
        \emph{conjectures} themselves
        (Conjectures~\ref{conj:T} and~\ref{conj:j}) involve
        infinite suprema and asymptotic decay, both of which are
        statable but presently unprovable in current Mathlib. The
        v3 Lean import goes as far as the spectral-radius bridge
        at a single $(T, j)$ and at a single depth $K_0$, but no
        further.
  \item An extension of the deterministic finite-residue-cell
        certificate to $K_0 > 16$ or to more refined residue
        $\mathrm{lift\_bits}$ values would tighten the bound but
        does not, by itself, change the scope: the certificate
        remains a bound on a finite substochastic retention
        matrix at a chosen depth, not a spectral-gap proof on an
        infinite operator.
\end{itemize}

\paragraph{Honest claim.}
We have reduced the open structural content of the Collatz
conjecture, at the level of the deterministic finite-residue-cell
$K_0 = 16$ taxonomy SCC, to a problem of \emph{spectral
perturbation theory on a finite explicit family of nonnegative
matrices}, with the foundational shadowing lemma, the finite
operator construction, the generic weighted Collatz--Wielandt
bound, and two specific generated finite numerical certificates
mechanically verified in Lean. We do not claim the conjecture
itself.

\subsection{Planned next steps}
\label{sec:planned}

The original Priority~A (phantom-set taxonomy) and Priority~B
(Lean formalization of the episode graph and transfer operator)
of v2 are completed in v3. The remaining priority is the analytic
program of Priority~C, now expressed as a research roadmap with
explicit gates rather than as a single outcome target.

\paragraph{Priority A (v2) --- Phantom-set taxonomy.\
Completed in v3.}
A finite enumeration of all primitive expansive compositions up to
depth $K_0 = 16$ is carried out in Section~\ref{sec:taxonomy}.
The taxonomy run exhibits a single nontrivial strongly connected
component of $1222$ raw states; the deterministic
finite-residue-cell construction on the compressed $(K, b)$
macro-state space yields an exact rational Collatz--Wielandt
bound $< 3/4$. Sensitivity at finer residue refinement
($\mathrm{lift\_bits} \in \{4, 5, 6\}$) shows the bound is
monotonically improving. The certificate is imported into Lean as
\ttbr{k16s16KDeterministicGeneratedSpectralRadiusBound}.

\paragraph{Priority B (v2) --- Lean formalization of the episode
graph and transfer operator.\ Completed in v3.}
The episode graph, the phase-state quotient, the truncated
transfer operator, the empirical signature decomposition, and the
generic weighted Collatz--Wielandt bound are formalized in
Section~\ref{sec:formalization} under the new modules
\texttt{EpisodeGraph.lean}, \texttt{Operator.lean}, and
\texttt{Bound.lean}. The single-$(T, j)$ certificate at
$T = 10, j = 32$ is generated and Lean-checked under
\texttt{Generated/T10J32HighBitTailCW.lean} with no
\texttt{axiom}, \texttt{sorry}, or \texttt{admit}.

\paragraph{Priority C --- Transfer-operator roadmap toward
Conjecture~\ref{conj:T}, gated.}
We do not pursue here an analytic Lasota--Yorke / Hennion /
Keller--Liverani argument. We instead lay out the program as a
sequence of producible artifacts, with an explicit fork to a
finite-rank fallback at the gate where the infinite phase model
is first interrogated. The list below mirrors the operational
roadmap in \texttt{lean/TODO.md}, Phase~10:
\begin{enumerate}
  \item \textbf{Literature and hypothesis matrix.} A table
        documenting, for each candidate framework
        (Keller--Liverani \cite{kellerliverani}, Hennion's
        compactness theorem, Baladi's positive transfer
        operator framework \cite{baladi}, Sarig's countable
        Markov shifts, GDMS systems, and the non-Archimedean
        thermodynamic formalism of \cite{anashin,bcm2020}),
        explicit theorem hypotheses against the present
        setting's known features, with compatibility against
        \cite[\S 7, Remark~9.17]{chang2026} as a column rather
        than a separate task.
  \item \textbf{Infinite phase space and quotient.} The single
        load-bearing question of the program: identify the
        infinite phase space whose finite
        $\full_{T, j}$ matrices are exact projections (or
        document that no such natural infinite kernel exists,
        in which case the program is to be terminated and the
        finite-rank certificate of step (G) is adopted as the
        program output).
  \item \textbf{Candidate Banach spaces.} For each candidate
        space (Lipschitz/H\"older on
        $\Zp \times \mathrm{Fin}\,4$; pesata variation on the
        residue tree; weighted countable shift), the strong
        and weak norms, the compactness or tightness
        mechanism, and the proof obligations.
  \item \textbf{Explicit transfer-operator formulation.}
        Orientation (push-forward vs.\ pull-back), weight,
        terminal killing, normalization, and the relation to
        the finite matrices.
  \item \textbf{Target Lasota--Yorke inequality.} Lemma
        statement with explicit constants or parameters and a
        checklist of local branch / distortion / tail estimates.
  \item \textbf{Conditional skeletons (Hennion + Keller--Liverani).}
        Conditional propositions producing essential
        spectral-radius bound under (E) and compactness, and
        perturbation control of the dominant eigenvalue along
        a truncation family.
  \item \textbf{Finite-rank / truncated-operator fallback,
        partial baseline.} The deterministic finite-residue-cell
        Collatz--Wielandt certificate of
        Section~\ref{sec:deterministic} is already such a
        baseline at $K_0 = 16, \mathrm{lift\_bits} = 4, 5, 6$;
        the program output here is its extension to $K_0 = 20$
        and to a refined-residue family, or its declaration as
        the program output if (B) concludes that no natural
        infinite kernel exists.
  \item \textbf{Computational experiments for constants.}
        Reproducible diagnostics for tails, branch regularity,
        mixed-norm drift, and CW-vector regularity, only after
        (C) and (D) fix the analytic framework.
  \item \textbf{Collaboration package.} A concise dossier with
        definitions, finite certificates, open lemmas, scripts,
        and exact questions, prepared to be readable by an
        external expert in transfer-operator analysis.
  \item \textbf{Decision and kill-switch.} A written
        go/no-go decision: v4 paper extending the present line,
        companion paper if the analytic content stabilises
        independently, or coauthored paper if the analytic
        layer is closed in collaboration. The kill-switch is
        binding: if step (B) concludes that the finite matrices
        are not exact projections of any natural infinite
        kernel, steps (E) and (F) are skipped and (G) becomes
        the program output.
\end{enumerate}
We do not have the analytic expertise to carry out steps
(E)--(F) in isolation, and we consider the present program a
collaboration target. The author is open to collaboration as
described in Section~\ref{sec:methodology}.

\section{Methodology and AI disclosure}
\label{sec:methodology}

This paper is the first artifact of a personal research program by
the author, an independent researcher, exploring a question of
methodology as much as one of mathematics:
\begin{quote}
\emph{To what extent can sustained collaboration with general-purpose
large language model assistants enable an independent researcher to
make non-standard, well-founded, externally reviewable attempts at
open mathematical problems?}
\end{quote}
The question is not a rhetorical one. It is the experiment of which
the present manuscript --- and the parallel Lean formalization
described in Section~\ref{sec:formalization} --- are concrete
artifacts. \textbf{Verifying what such collaboration can and cannot
deliver, in a setting where each step is reproducible and openly
reviewable, is an explicit secondary goal of this work}, on equal
footing with the substantive Collatz progress.

\subsection{Why this disclosure matters}

A clear separation of contributions is essential for the reader's
correct evaluation of the material. Errors that would normally be
caught by departmental colleagues, weekly seminars, or referee
comments are not part of the development cycle here; they have to
be substituted by other safeguards (numerical reproducibility,
mechanical verification, multi-AI cross-checking, public release
for external scrutiny).

\subsection{Author's contribution}

The conceptual framing is due to the author: the gravitational
debt metaphor for orbit expansion, the empirical observation that
$\nu_2 = 1$ corridors are the locus of explosive growth, the
structural intuition that integer rebels shadow $2$-adic phantom
orbits associated with negative rational fixed points of the
Syracuse map, the strategic pivot from the empirical
quasi-Lyapunov candidates of the early phase to the present
spectral program (after their false-alarm pattern proved
structurally inseparable from genuine rebels), and the decision to
publish at this stage. The author led the research direction
throughout: the choices of what to investigate, what to refine,
when to stop scaling, and what to publish are the author's. The
author is also responsible for the literature reconnaissance
phase that produced the related-work analysis of
Section~\ref{sec:related}.

\subsection{AI assistance, by phase}

\paragraph{Early empirical phase (\texttt{scripts/early\_empirical/}).}
Scripts numbered roughly $01$ through $48$, plus
\texttt{validate\_*.py}. Topics: orbit profiling, debt
accumulation metrics, family identification in inverse trees,
false-alarm classification, empirical search for the candidate
quasi-Lyapunov $H$ function. Developed with assistance from
\textbf{OpenAI Codex} and \textbf{Google Gemini}. The author
guided the empirical questions and validated the results manually
against worked examples.

\paragraph{Spectral program phase (\texttt{scripts/spectral\_program/}).}
Scripts numbered $49$ through $87$. Topics: phantom rational
representatives, shadowing congruence verification, episode graph
construction, strongly connected component identification, the
critical symbolic operator, the CORE/TAIL decomposition, the
weighted Collatz--Wielandt bound, truncation stability, per-node
and cross-node SCC certificates. These scripts, the mathematical
formalization of Lemma~\ref{lem:shadowing} and the weighted
Collatz--Wielandt bound, the construction of the transfer
operator, the conditional reduction theorem, and the entire
\LaTeX{} manuscript were developed with substantial assistance
from \textbf{Anthropic's Claude}, accessed via the Claude Code
interface in extended interactive sessions.

In particular, the following content originated in conversation
with Claude:
\begin{itemize}
  \item the proof of Lemma~\ref{lem:shadowing} (induction on the
        bit valuation of the difference $S^j(n) - S^j(q_w)$);
  \item the construction of the refined phase state
        $(\nu_2(t), \mathrm{odd}(t) \bmod 4, h \bmod 4)$;
  \item the empirical core/tail signature decomposition;
  \item the derivation of the weighted Collatz--Wielandt bound and
        the observation that the right Perron eigenvector of
        $\full$ (rather than $\core$) is the correct test vector;
  \item the cross-node operator construction and the topological
        interpretation of the SCC as a star centered on
        $(12, 2, 1)$;
  \item the formulation of the open conjectures and the conditional
        reduction theorem.
\end{itemize}
The author's role in this phase was to ask follow-up questions, to
suggest experiments, to validate against numerical sanity checks,
and to make strategic decisions about scope and publication
timing.

\paragraph{Lean formalization phase, v1$\to$v2
(\texttt{lean/}).}
The Lean~4 + Mathlib formalization of Lemma~\ref{lem:shadowing}
and Corollary~\ref{cor:no-infinite}
(Section~\ref{sec:formalization}) was carried out in a sequence
of focused sessions across approximately three calendar days,
alternating \textbf{Claude} (Claude Code interface) and
\textbf{OpenAI Codex} as primary technical assistants. Cross-AI
checking was used to identify mistakes that one assistant did
not catch: in particular, several Mathlib API mismatches arising
from the rapid evolution of Mathlib lemma names (e.g.\
deprecated \texttt{push\_neg}, the maximal-ideal API renaming,
\texttt{List.get?} removal in v4.29.1) were identified by one
assistant and fixed by the other. The author's role was to set
session-level goals, validate that each session's claimed
artifact was honest (zero \texttt{sorry}, \texttt{lake build}
clean, statements faithful to the paper's language), and to
maintain the operational \texttt{TODO.md} across sessions. The
session log at the bottom of \texttt{lean/TODO.md} records the
exact sequence of steps and which assistant produced which
artifact.

\paragraph{Phantom-set taxonomy phase, v2$\to$v3
(\texttt{scripts/phantom\_taxonomy/}).}
The taxonomy construction of Section~\ref{sec:taxonomy} was
carried out as a sequence of self-contained Python scripts:
necklace enumeration, exact rational phantom representatives,
orbit harness, sampled SCC integration, and the final
deterministic finite-residue-cell generator. The Lean import was
mediated by two generator scripts
(\texttt{lean\_scc\_certificate.py},
\texttt{lean\_cw\_summary.py}) that emit
\texttt{Generated/*.lean} modules from the JSON certificates. The
deterministic generator was developed in adversarial review
against the empirical sampled generator: any divergence between
the two on a shared row of the macro-state matrix was flagged
for inspection, and the closure of the deterministic
$\mathrm{lift\_bits} = 4$ enumeration to zero budget exits is
the integrity check that both generators agree on the support of
the retention matrix.

\paragraph{Lean formalization phase, v2$\to$v3
(\texttt{lean/CollatzShadowing/}).}
The Phase~8 modules (\texttt{EpisodeGraph.lean},
\texttt{Operator.lean}, \texttt{Bound.lean}) and the Phase~8
generated certificates (\texttt{Generated/T10CriticalSymbolic},
\texttt{T10J32HighBitTail}, \texttt{T10J32HighBitTailCW},
\texttt{K16S16KSCC}, \texttt{K16S16KExactCWSummary},
\texttt{K16S16KBridge}, \texttt{K16S16KDeterministicCW}) were
developed in extended interactive sessions with
\textbf{Claude} and \textbf{OpenAI Codex}, with the latter
typically driving the generator scripts and the former driving
the supporting Mathlib API design and the spectral-radius
bridge. A central methodological observation from this phase is
the removal of the residual \texttt{axiom} that the early
$T = 10, j = 32$ CW certificate had used to bridge the
\texttt{decide}-style numerical check to the
\texttt{spectralRadius} target: the final
\texttt{Generated/T10J32HighBitTailCW.lean} closes the bridge
without \texttt{axiom}, \texttt{sorry}, or \texttt{admit}, with
the 224-row certificate split across 14 row modules to keep
elaboration tractable. The cross-AI protocol was used here
exactly as in v1$\to$v2: when one assistant introduced a
\texttt{sorry} or an \texttt{axiom}, the other was asked to
either close it or to explain why it could not be closed
without enlarging the scope.

\paragraph{Literature reconnaissance phase
(\texttt{arxiv-affini/}, separate from the project repository).}
The related-work content of Section~\ref{sec:related} is the
output of a separate AI-assisted reconnaissance: a corpus of
$\sim 100$ arXiv papers in math.NT, math.DS, math.SP, math.FA, and
related categories was assembled, ranked by topical proximity to
the present work, and reviewed for explicit overlaps and
differences. The author validated the most consequential
candidates (Chang \cite{chang2026}, Siegel \cite{siegel2023a},
Rozier \cite{rozier2025}) by direct inspection of the relevant
sections, not only via abstracts. The reconnaissance corpus and
its automated ranking scripts are not committed to the public
project repository: they consist largely of third-party PDFs and
of intermediate scratch documents whose contents are subsumed by
Section~\ref{sec:related} of this paper. The methodology of the
reconnaissance is, however, openly documented; see the
``Literature reconnaissance method'' note appended to
\texttt{METHODOLOGY.md}.

\subsection{Cross-AI verification protocol}

The methodology that emerged in practice, and which we recommend
for analogous attempts, is a four-stage verification cycle for
every non-trivial mathematical claim:
\begin{enumerate}
  \item the claim is first developed in dialogue with a primary
        AI assistant;
  \item it is then re-examined, in cold-start form, by a second
        AI assistant whose role is adversarial review;
  \item where applicable, it is reduced to a deterministic
        numerical experiment whose output can be inspected
        against worked examples by the author;
  \item where applicable, it is mechanically verified in Lean.
\end{enumerate}
Stages~3 and~4 are the strongest safeguards: a numerical
experiment that fails to reproduce a claimed identity exposes the
error immediately, and a Lean proof that fails to type-check
exposes any logical gap that the dialogue rounded over.

\subsection{Verification status, current}

\begin{itemize}
  \item All numerical results in Sections~\ref{sec:numerics},
        \ref{sec:cross}, and \ref{sec:taxonomy} are deterministic
        and reproducible from the supplementary scripts; anyone
        with Python 3.10+, NumPy, and SciPy can re-run them
        (including the exact-rational deterministic
        finite-residue-cell certificate of
        Section~\ref{sec:deterministic}).
  \item Lemma~\ref{lem:shadowing},
        Corollary~\ref{cor:no-infinite}, the generic
        weighted Collatz--Wielandt bound, the
        $T = 10, j = 32$ single-node certificate
        $\spec(\full_{10,32}) \le 97/2000$, and the $K_0 = 16$
        deterministic finite-residue-cell certificate
        $\spec(P_{K, b}) \le 3/4$ are mechanically verified
        in Lean~4 with Mathlib; the project is
        \texttt{sorry}-free and \texttt{axiom}-free. See
        Section~\ref{sec:formalization} and \texttt{lean/}.
  \item As of this version, the work has not been reviewed by a
        human mathematical expert. The construction of the
        transfer operator, the weighted bound, the conditional
        reduction theorem, and the taxonomy certificate rest on
        AI-assisted formalization that has not yet been
        externally validated. External expert review is invited.
  \item The literature reconnaissance is necessarily limited by
        the corpus consulted ($\sim 100$ papers up to and
        including early 2026). Between v2 (2026-05-10) and v3
        (2026-05-13) no new arXiv work in the directly
        overlapping niche surfaced; the standing monthly arXiv
        scan described in
        Section~\ref{sec:methodology-future-search} continues to
        apply.
\end{itemize}

\subsection{Future search for analogous studies}
\label{sec:methodology-future-search}

The risk of inadvertent duplication is highest in the
intersection of (i) symbolic / $p$-adic dynamics, (ii) finite
transfer operators, (iii) Collatz-like iterations. The discovery
of \cite{chang2026} two months after our v1 was deposited is a
concrete instance of this risk. To reduce it for future iterations
of this work and for adjacent problems, the author commits to the
following recurring search procedure (also documented in
\texttt{METHODOLOGY.md}):
\begin{enumerate}
  \item monthly arXiv listing scans in \texttt{math.NT},
        \texttt{math.DS}, \texttt{math.SP}, \texttt{math.FA},
        \texttt{math.CO}, with manual screening of titles and
        abstracts for the keywords \emph{phantom}, \emph{shadow},
        \emph{transfer operator}, \emph{Collatz}, \emph{Syracuse},
        \emph{$3x+1$}, \emph{2-adic}, \emph{spectral gap},
        \emph{Collatz--Wielandt}, \emph{Perron};
  \item targeted Google Scholar and Semantic Scholar queries with
        the same keyword combinations, plus citation-graph
        traversal of the seed papers \cite{chang2026,
        siegel2023a, rozier2025, neklyudov2022,
        leventides2020, lemmensnussbaum};
  \item AMS MathSciNet and zbMATH searches under MSC primary
        $11B83$ and $37N99$, with secondary $47A10$ and $11S99$;
  \item monitoring of the arXiv 2603.* and 2604.* releases that
        cite the seed papers above;
  \item should any future related work surface, an explicit
        addendum or revised version note will be prepared and
        added to the public repository within a reasonable
        time-frame.
\end{enumerate}
The author considers this duty of continuous reconnaissance an
intrinsic part of an AI-assisted research program: standard
in-person community signals (department seminars, ad-hoc
discussions, workshop announcements) that ordinarily flag related
work for institutionally embedded researchers are absent here, so
the surveillance has to be made explicit.

\subsection{How to engage}
\begin{itemize}
  \item \textbf{For the math.} Read Sections~\ref{sec:shadowing},
        \ref{sec:operator}, \ref{sec:bound}, and
        \ref{sec:taxonomy} of this paper. If you find an error,
        an unjustified step, or an implicit assumption worth
        flagging, please contact the author at the address on the
        title page.
  \item \textbf{For the Lean formalization.} Clone the repository
        and run \texttt{lake build} from the \texttt{lean/}
        directory. The headline v3 declarations
        (\texttt{exact\_shadowing},
        \texttt{exact\_shadowing\_periods},
        \texttt{no\_infinite\_period\_congruence\_expansive},
        \ttbr{t10j32HighBitTailSpectralRadiusBound\_97\_2000},
        \ttbr{k16s16KDeterministicGeneratedSpectralRadiusBound})
        compile without \texttt{sorry} and without \texttt{axiom}
        on Lean v4.29.1 and Mathlib v4.29.1.
  \item \textbf{For the numerics.} Clone the repository, install
        NumPy and SciPy, and re-run any of the scripts in
        \texttt{scripts/spectral\_program/} or
        \texttt{scripts/phantom\_taxonomy/}. The deterministic
        finite-residue-cell certificate can be regenerated and
        re-verified with the single command sequence documented
        in
        \texttt{notes/phantom\_taxonomy\_empirical\_scc\_integration.md}.
  \item \textbf{For collaboration.} The author is open to
        collaboration with anyone willing to take the proofs
        seriously, and especially welcomes (a) attempts to make
        the gated Phase~10 roadmap of
        Section~\ref{sec:planned} concrete in the analytic
        framework of Keller--Liverani / Hennion / Lasota--Yorke
        on a Banach space adapted to the refined phase quotient
        (steps (A)--(F) of the roadmap), and (b) independent
        verification of the deterministic finite-residue-cell
        certificate of Section~\ref{sec:deterministic} at
        $K_0 = 16$ or its extension to $K_0 = 20$.
\end{itemize}

\begin{thebibliography}{99}

\bibitem{tao2019}
T.~Tao,
\emph{Almost all orbits of the Collatz map attain almost bounded values},
arXiv:1909.03562, 2019; published in Forum of Mathematics, Pi, 2022.

\bibitem{lagarias1985}
J.~C.~Lagarias,
\emph{The 3x+1 problem and its generalizations},
American Mathematical Monthly 92 (1985), 3--23.

\bibitem{lagarias2010}
J.~C.~Lagarias (ed.),
\emph{The Ultimate Challenge: The 3x+1 Problem},
American Mathematical Society, 2010.

\bibitem{laga2003bib}
J.~C.~Lagarias,
\emph{The 3x+1 Problem: An Annotated Bibliography (1963--1999)},
arXiv:math/0309224.

\bibitem{laga2009bib}
J.~C.~Lagarias,
\emph{The 3x+1 Problem: An Annotated Bibliography, II (2000--2009)},
arXiv:math/0608208.

\bibitem{kellerliverani}
G.~Keller and C.~Liverani,
\emph{Stability of the spectrum for transfer operators},
Annali della Scuola Normale Superiore di Pisa, Cl. Sci. (4) 28 (1999), 141--152.

\bibitem{baladi}
V.~Baladi,
\emph{Positive Transfer Operators and Decay of Correlations},
World Scientific, 2000.

\bibitem{anashin}
V.~Anashin,
\emph{The non-Archimedean theory of discrete systems},
Math. Comput. Sci. 6 (2012), 375--393.

\bibitem{chang2026}
E.~Y.~Chang,
\emph{Exploring Collatz Dynamics with Human-LLM Collaboration},
arXiv:2603.11066, 2026.

\bibitem{siegel2023a}
M.~C.~Siegel,
\emph{The Collatz Conjecture and Non-Archimedean Spectral Theory --- Part I:
Arithmetic Dynamical Systems and Non-Archimedean Value Distribution Theory},
arXiv:2007.15936, 2020--2023.

\bibitem{siegel2023b}
M.~C.~Siegel,
\emph{The Collatz Conjecture and Non-Archimedean Spectral Theory --- Part I.5:
How To Write The (Weak) Collatz Conjecture As A Contour Integral},
arXiv:2111.07883, 2021--2023.

\bibitem{siegel2026hydra}
M.~C.~Siegel,
\emph{The Hydra Map and Numen Formalisms for Collatz-Type Problems},
arXiv:2601.17030, 2026.

\bibitem{rozier2025}
O.~Rozier,
\emph{Parity sequences of the 3x+1 map on the 2-adic integers and Euclidean embedding},
arXiv:1805.00133, 2018--2025.

\bibitem{neklyudov2022}
M.~Neklyudov,
\emph{Functional analysis approach to the Collatz conjecture},
arXiv:2106.11859, 2022.

\bibitem{berg1996}
L.~Berg and G.~Meinardus,
\emph{Functional equations connected with the Collatz problem},
Results in Mathematics 25 (1994), 1--12; further work 1990s.

\bibitem{leventides2020}
J.~Leventides and C.~Poulios,
\emph{Koopman operators and the 3x+1-dynamical system},
arXiv:2010.12987, 2020.

\bibitem{mori2025}
T.~Mori,
\emph{Application of Operator Theory for the Collatz Conjecture},
arXiv:2411.08084, 2025.

\bibitem{laarhoven2012}
T.~Laarhoven and B.~de~Weger,
\emph{The Collatz conjecture and De Bruijn graphs},
arXiv:1209.3495, 2012.

\bibitem{lemmensnussbaum}
M.~Akian, S.~Gaubert, and R.~Nussbaum,
\emph{A Collatz--Wielandt characterization of the spectral radius
of order-preserving homogeneous maps on cones},
arXiv:1112.5968, 2011--2014.

\bibitem{akiangaubert}
M.~Akian and S.~Gaubert,
\emph{Spectral theorem for convex monotone homogeneous maps,
and ergodic control},
Nonlinear Anal. 52 (2003), 637--679.

\bibitem{bcm2020}
J.~Bastos, D.~Caprio, and A.~Messaoudi,
\emph{Shadowing and Stability in p-adic dynamics},
arXiv:2001.02737, 2020.

\bibitem{quasicompact}
\emph{Quasi-compactness of Frobenius--Perron Operator for
Piecewise Convex Maps with Countable Branches},
arXiv:2406.19929, 2024.

\bibitem{certifiedspectral}
\emph{Certified spectral approximation of transfer operators
and the Gauss map},
arXiv:2602.19435, 2026.

\bibitem{inselmann2024}
M.~Inselmann,
\emph{An approximation of the Collatz map and a lower bound
for the average total stopping time},
arXiv:2402.03276, 2024.

\end{thebibliography}

\appendix

\section{Phantom records used in this work}
\label{app:phantoms}

The expansive phantom representatives identified empirically up to
$k = 24$ and used as monitored nodes in Section~\ref{sec:episodes}:

\begin{center}
\begin{tabular}{rl}
\toprule
$k$ & $q_w \approx$ \\
\midrule
10 & $-5\,739\,192\,681\,529 / 2\,404\,426\,874\,857 \approx -2.387$ \\
11 & $-2\,199\,012\,694\,031 / 709\,849\,655\,971 \approx -3.098$ \\
12 (cycle 1) & $-2123 / 1675 \approx -1.267$ \\
12 (cycle 2) & $-817 / 601 \approx -1.359$ \\
20 & $-104\,312\,644\,103 / 30\,307\,317\,785 \approx -3.442$ \\
\bottomrule
\end{tabular}
\end{center}

The full list and the construction of these representatives from
the rational fixed point formula \eqref{eq:qw} is implemented in
\verb|54_phantom_rational_shadows.py|. By
\cite[Theorem~7.15]{chang2026} this list is a small subset of the
infinite phantom universe; we use it as a finite working set for
the empirical SCC analysis of Section~\ref{sec:cross}.

\section{Supplementary scripts and Lean code}
\label{app:scripts}

All supporting material is available in the GitHub repository
\begin{center}
\url{https://github.com/PieroBorgatta/Collatz}
\end{center}
The principal Python scripts and what each computes:

\begin{description}[leftmargin=4em, style=nextline]
  \item[\texttt{54\_phantom\_rational\_shadows.py}]
    Computes the $2$-adic representatives of expansive phantom
    cycles; underpins all subsequent constructions.
  \item[\texttt{55\_shadowing\_congruence.py}]
    Verifies Lemma~\ref{lem:shadowing} on the empirical phantom
    set for $b \le 10$.
  \item[\texttt{59\_shadow\_episode\_graph.py},\
        \texttt{60\_episode\_graph\_diagnostics.py}]
    Constructs the episode graph and identifies its SCCs.
  \item[\texttt{75\_critical\_symbolic\_operator.py}]
    Builds the truncated finite operator at the critical node.
  \item[\texttt{77\_high\_bit\_tail\_bound.py}]
    Decomposes $\full = \core + \tail$ from the empirical
    signature majority.
  \item[\texttt{82\_boundary\_scc\_spectral\_scaling.py}]
    Spectral scaling on SCCs touching the boundary layer
    $\nu_2(t) \ge T$.
  \item[\texttt{83\_perturbation\_gap\_scaling.py}]
    Measures the gap $\spec(\full) - \spec(\core)$.
  \item[\texttt{84\_weighted\_perturbation\_bound.py}]
    Computes the weighted bound \eqref{eq:bound}.
  \item[\texttt{85\_truncation\_stability.py}]
    Cauchy/signature stability of the truncation in $j$.
  \item[\texttt{86\_scc\_node\_certificates.py}]
    Per-node bounds on the SCC.
  \item[\texttt{87\_cross\_node\_transfer.py}]
    Assembled cross-node transfer operator on the SCC.
\end{description}

The v3 phantom-taxonomy scripts are in
\texttt{scripts/phantom\_taxonomy/}:

\begin{description}[leftmargin=4em, style=nextline]
  \item[\texttt{necklace\_counts.py}]
    M\"obius enumeration of primitive cyclic compositions
    $M(K, \ell)$ with brute-force verification through
    $K \le 10$ (Section~\ref{sec:necklace}).
  \item[\texttt{phantom\_representatives.py}]
    Exact-rational $2$-adic representatives for every primitive
    expansive composition; CSV outputs
    \texttt{phantom\_representatives\_k3\_16.csv} and
    \texttt{phantom\_representatives\_k3\_20.csv}
    (Section~\ref{sec:representatives}).
  \item[\texttt{orbit\_harness.py}]
    Integer-lift orbit harness that monitors phantom-congruence
    hits and emits per-source CSVs of nodes, events, and edges
    (Section~\ref{sec:harness}).
  \item[\texttt{scc\_transfer\_summary.py},\
        \texttt{scc\_collatz\_wielandt.py},\
        \texttt{scc\_cw\_certificate.py}]
    Substochastic-retention summary, empirical
    Collatz--Wielandt expression, and exact rational
    verification of CW certificates
    (Section~\ref{sec:empirical-baseline}).
  \item[\texttt{deterministic\_residue\_transfer.py}]
    Deterministic enumeration of all $2^{\mathrm{lift\_bits}}$
    finite residue subclasses per source state, producing exact
    rational deterministic matrices plus coverage manifest
    (Section~\ref{sec:deterministic}).
  \item[\texttt{lean\_scc\_certificate.py},\
        \texttt{lean\_cw\_summary.py},\
        \texttt{lean\_phase\_transfer.py},\
        \texttt{lean\_high\_bit\_tail.py},\
        \texttt{lean\_t10j32\_cw.py},\
        \texttt{lean\_cw\_smoke.py},\
        \texttt{export\_high\_bit\_tail\_edges.py}]
    Generators emitting the v3
    \texttt{CollatzShadowing.Generated.*.lean} modules
    (Section~\ref{sec:deterministic-lean} and
    Section~\ref{sec:formalization}).
\end{description}

The Lean~4 + Mathlib formalization is in the subdirectory
\texttt{lean/} of the repository:

\begin{description}[leftmargin=4em, style=nextline]
  \item[\texttt{lean/CollatzShadowing/Basic.lean}]
    The accelerated Syracuse map $S\colon \N \to \N$ and the
    valuations $\nu_2$ on $\N$, $\Z$, $\Q$, $\Zp$.
  \item[\texttt{lean/CollatzShadowing/Phantom.lean}]
    The \texttt{PhantomWord} datatype, affine fold of the
    Syracuse step, the named coefficients $C_w$, $A_w$, the
    rational fixed point $q_w$, and its $\Zp$ representative.
  \item[\texttt{lean/CollatzShadowing/Syracuse2Adic.lean}]
    The total $2$-adic extension $\mathtt{Syracuse2adic}\colon
    \Zp \to \Zp$ and the bridge to $S$ on $\N$.
  \item[\texttt{lean/CollatzShadowing/Auxiliary.lean}]
    Supporting lemmas: $\nu_2(3 n) = \nu_2(n)$ on $\Zp$,
    odd-denominator of $q_w$, the affine difference identity
    $S^j(n) - S^j(q_w) = (3^j / 2^{B_j})(n - q_w)$, the strict
    ultrametric stability of $\nu_2$ under proximity, the
    matching property of $q_w$'s own orbit, the periodicity
    identity $B(b L) = b A$, and the closed-form identity
    \texttt{B\_closed\_form}.
  \item[\texttt{lean/CollatzShadowing/Shadowing.lean}]
    Lemma~\ref{lem:shadowing}
    (\texttt{exact\_shadowing}) and its periodic specialization
    (\texttt{exact\_shadowing\_periods}); both
    \texttt{sorry}-free.
  \item[\texttt{lean/CollatzShadowing/NoInfinite.lean}]
    Corollary~\ref{cor:no-infinite}
    (\texttt{no\_infinite\_period\_congruence\_expansive});
    \texttt{sorry}-free.
  \item[\texttt{lean/CollatzShadowing/EpisodeGraph.lean}]
    \texttt{EpisodeNode}, the finite cutoff box
    \texttt{TruncatedEpisodeNode K C B}, the edge relation
    \texttt{EpisodeGraph.edge}, reachability and SCC wrappers,
    and the concrete finite-walk certificate format
    \texttt{HubSCCCertificate} (v3 addition).
  \item[\texttt{lean/CollatzShadowing/Operator.lean}]
    The phase-state quotient
    \texttt{PhaseState V := Fin (V+1) $\times$ Fin 4 $\times$ Fin 4},
    the operator interfaces \texttt{TransferMatrix V},
    \texttt{ProbabilityEntry}, \texttt{RowSubstochastic}, and
    the bundled \texttt{OperatorDecomposition} (v3 addition).
  \item[\texttt{lean/CollatzShadowing/Bound.lean}]
    The generic weighted Collatz--Wielandt finite-matrix bound
    (\texttt{FiniteCWCertificate},
    \ttbr{spectralRadius\_le\_of\_finiteCWCertificate},
    \texttt{OperatorDecomposition.spectralRadius\_le\_full})
    (v3 addition).
  \item[\texttt{lean/CollatzShadowing/Generated/T10CriticalSymbolic.lean},\
        \texttt{T10J32HighBitTail.lean},\
        \texttt{T10J32HighBitTailCW.lean},\
        \texttt{T10J32HighBitTailCWData.lean},\
        \texttt{T10J32HighBitTailCWRows00.lean--Rows13.lean}]
    The $T = 10, j = 32$ generated finite-operator
    Collatz--Wielandt certificate and its spectral-radius bridge
    (v3 addition).
  \item[\texttt{lean/CollatzShadowing/Generated/K16S16KSCC.lean},\
        \texttt{K16S16KExactCWSummary.lean},\
        \texttt{K16S16KBridge.lean}]
    The $K_0 = 16$ taxonomy SCC, the empirical 37-node CW
    certificate summary, and the bridge to spectral-radius
    (v3 addition; historical baseline).
  \item[\texttt{lean/CollatzShadowing/Generated/K16S16KDeterministicCW.lean}]
    The deterministic finite-residue-cell CW certificate at
    $K_0 = 16, \mathrm{lift\_bits} = 4$ and its spectral-radius
    bridge
    \ttbr{k16s16KDeterministicGeneratedSpectralRadiusBound}
    (v3 addition; the v3 headline result on the taxonomy SCC).
  \item[\texttt{lean/CollatzShadowing/Generated/K16S16KCWSmoke.lean}]
    A single-row smoke certificate kept as a tractable
    regression test for the CW pipeline.
  \item[\texttt{lean/CollatzShadowing/EPISODE\_INVENTORY.md}]
    Mathlib API inventory for the Phase-8 modules
    (v3 addition).
  \item[\texttt{lean/CollatzShadowing/INVENTORY.md}]
    A complete inventory of the Mathlib API used and of the
    chosen modeling decisions (in particular, congruence as
    ideal membership in $\Zp$).
\end{description}
Building requires Lean v4.29.1 and Mathlib v4.29.1, both pinned
in the \texttt{lean-toolchain} and \texttt{lakefile.toml} files.

\end{document}
